% संपूर्ण मराठी ग्रंथातील प्रकरण 10: नैसर्गिक निगमन
% हा संपादनयोग्य स्रोतखंड आहे. संकलनासाठी संपूर्ण स्रोतसंग्रहातील
% अक्षररूपे, पूर्वभाग, चिन्हव्याख्या आणि आकृती-स्रोत आवश्यक आहेत.
% मूळ: Open Logic Project; मजकूर आणि रूपांतर CC BY 4.0.
% यंत्रानुवाद, दुरुस्ती आणि तपासणी: OpenAI Codex —
% GPT-5.6 Sol आणि GPT-6 Sol, दोन्ही Ultra effort.
% दुरुस्ती-पुनरावलोकन: GPT-6.1 Sol, Ultra effort.
\chapter{नैसर्गिक निगमन}\label{fol:ntd::chap}
\begin{editorial}
हे प्रकरण Gentzen/Prawitz शैलीतील नैसर्गिक निगमन पद्धत मांडते.

नैसर्गिक निगमनाशी सिद्धता-पद्धती म्हणून संबंधित सामग्री समाविष्ट किंवा
वगळण्यासाठी ``prfND'' हा टॅग वापरा.
\end{editorial}






\section{नियम आणि निष्पत्ती}\label{fol:ntd:rul:sec}

\begin{explain}
नैसर्गिक निगमन पद्धतींचा उद्देश गणितीय सिद्धतेत वापरल्या जाणाऱ्या अनौपचारिक
तर्कपद्धतीशी जवळचे साम्य राखणे हा आहे (म्हणूनच ती काहीशी ``नैसर्गिक'' आहे).
नैसर्गिक निगमनातील सिद्धता गृहीतकांपासून सुरू होतात.
त्यानंतर अनुमाननियम लागू केले जातात. \Intro{\lnot}, \Intro{\lif},
\Elim{\lor} आणि \Elim{\lexists} अनुमाननियमांद्वारे
गृहीतके ``मुक्त'' केली जातात; स्पष्टतेसाठी मुक्त केलेल्या
गृहीतकाचा लेबल त्या अनुमानाजवळ ठेवला जातो.
\end{explain}

\begin{defn}[गृहीतक]
\emph{गृहीतक} म्हणजे कोणत्याही शाखेच्या सर्वांत वरच्या स्थानावर असलेले
कोणतेही वाक्य होय.
\end{defn}

नैसर्गिक निगमनातील निष्पत्त्या ही वाक्यांची विशिष्ट वृक्षरचना असते;
त्यातील सर्वांत वरच्या वाक्ये गृहीतके असतात. एखादे वाक्य एक,
दोन किंवा तीन इतर क्रमवर्तींच्या खाली असल्यास, ते अनुमाननियमाने योग्य
प्रकारे आलेले असले पाहिजे. अनुमानाच्या शीर्षस्थानी असलेल्या वाक्यांना
\emph{आधारविधाने} म्हणतात आणि खालील वाक्याला अनुमानाचा
\emph{निष्कर्ष} म्हणतात. प्रत्येक संकारक साठी नियम जोड्यांमध्ये येतात:
एक प्रवेशन नियम आणि एक विलोपन नियम. ते निष्कर्षात संकारक आणतात
किंवा नियमाच्या आधारविधानातून संकारक काढतात. काही नियम विशिष्ट
प्रकारचे गृहीतक \emph{मुक्त} करण्याची परवानगी देतात. कोणते
गृहीतक कोणत्या अनुमानाने मुक्त केले हे दर्शवण्यासाठी गृहीतक आणि
अनुमान दोन्हींना लेबल देतो. हे गृहीतक ``$\Discharge{\varphi}{n}$.` असे लिहून दर्शवले जाते.


सर्व संकारक $\land$, $\lor$, $\lif$, $\lnot$ आणि $\lfalse$ यांच्यासाठीचे नियम विचारात घेणे प्रचलित आहे, जरी त्यांपैकी काही परिभाषित केलेले असले तरी.











\section{विधानीय नियम}\label{fol:ntd:prl:sec}


\subsection{$\land$ साठीचे नियम}

\begin{defish}
\AxiomC{$\varphi$}
\AxiomC{$\psi$}
\RightLabel{\Intro{\land}}
\BinaryInfC{$\varphi \land \psi$}
\DisplayProof
\hfill
\begin{tabular}{r}
\AxiomC{$\varphi \land \psi$}
\RightLabel{\Elim{\land}}
\UnaryInfC{$\varphi$}
\DisplayProof
\\[3ex]
\AxiomC{$\varphi \land \psi$}
\RightLabel{\Elim{\land}}
\UnaryInfC{$\psi$}
\DisplayProof
\end{tabular}
\end{defish}
\subsection{$\lor$ साठीचे नियम}
\begin{defish}
\begin{tabular}{r}
\AxiomC{$\varphi$}
\RightLabel{\Intro{\lor}}
\UnaryInfC{$\varphi \lor \psi$}
\DisplayProof
\\[3ex]
\AxiomC{$\psi$}
\RightLabel{\Intro{\lor}}
\UnaryInfC{$\varphi \lor \psi$}
\DisplayProof
\end{tabular}
\hfill
\AxiomC{$\varphi \lor \psi$}
\AxiomC{$\Discharge{\varphi}{n}$}
\DeduceC{$\chi$}
\AxiomC{$\Discharge{\psi}{n}$}
\DeduceC{$\chi$}
\DischargeRule{\Elim{\lor}}{n}
\TrinaryInfC{$\chi$}
\DisplayProof
\end{defish}
\subsection{$\lif$ साठीचे नियम}
\begin{defish}
\AxiomC{$\Discharge{\varphi}{n}$}
\DeduceC{$\psi$}
\DischargeRule{\Intro{\lif}}{n}
\UnaryInfC{$\varphi \lif \psi$}
\DisplayProof
\hfill
\AxiomC{$\varphi \lif \psi$}
\AxiomC{$\varphi$}
\RightLabel{\Elim{\lif}}
\BinaryInfC{$\psi$}
\DisplayProof
\end{defish}
\subsection{$\lnot$ साठीचे नियम}
\begin{defish}
\AxiomC{$\Discharge{\varphi}{n}$}
\noLine
\DeduceC{$\lfalse$}
\DischargeRule{\Intro{\lnot}}{n}
\UnaryInfC{$\lnot \varphi$}
\DisplayProof
\hfill
\AxiomC{$\lnot \varphi$}
\AxiomC{$\varphi$}
\RightLabel{\Elim{\lnot}}
\BinaryInfC{$\lfalse$}
\DisplayProof
\end{defish}
\subsection{$\lfalse$ साठीचे नियम}
\begin{defish}
\AxiomC{$\lfalse$}
\RightLabel{\FalseInt}
\UnaryInfC{$\varphi$}
\DisplayProof
\hfill
\AxiomC{$\Discharge{\lnot \varphi}{n}$}
\DeduceC{$\lfalse$}
\DischargeRule{\FalseCl}{n}
\UnaryInfC{$\varphi$}
\DisplayProof
\end{defish}
लक्षात घ्या की $\Intro{\lnot}$ आणि $\FalseCl$ हे अगदी सारखे आहेत: फरक
असा की $\Intro{\lnot}$ हा नकारयुक्त वाक्य~$\lnot \varphi$ निष्पन्न
करतो, तर $\FalseCl$ हा सकारात्मक वाक्य~$\varphi$ निष्पन्न करतो.
जेव्हा एखादा नियम एखादे गृहीतक मुक्त करता येते असे दर्शवतो, तेव्हा आपण ती
परवानगी मानतो, आवश्यकता नाही. उदा., $\Intro{\lif}$ नियमात आधारविधान~$\psi$
च्या निष्पत्तीमधील $\varphi$ या रूपाची कितीही गृहीतके—शून्यसुद्धा—आपण
मुक्त करू शकतो.
\section{संख्यापकांचे नियम}\label{fol:ntd:qrl:sec}
\subsection{$\lforall$ साठीचे नियम}
\begin{defish}
\AxiomC{$\varphi(a)$}
\RightLabel{\Intro{\lforall}}
\UnaryInfC{$\lforall[x][\Atom{\varphi}{x}]$}
\DisplayProof
\hfill
\AxiomC{$\lforall[x][\Atom{\varphi}{x}]$}
\RightLabel{\Elim{\lforall}}
\UnaryInfC{$\varphi(t)$}
\DisplayProof
\end{defish}
$\lforall$ च्या नियमांमध्ये, $t$ हे बंद पद आहे (म्हणजे, कोणतेही चर
नसलेले पद), आणि $a$~हा असा स्थिरांक आहे की तो
निष्कर्ष~$\lforall[x][\varphi(x)]$ मध्ये किंवा आधारविधान~$\varphi(a)$ ने समाप्त
होणाऱ्या निष्पत्तीमध्ये मुक्त न केलेली असलेल्या कोणत्याही
गृहीतकात येत नाही. आपण $a$ ला \Intro{\lforall} अनुमानाचा
\emph{आयगेन चर} म्हणतो.\footnote{वरील नियमात $a$ हा अचर असला तरी आपण
``आयगेन चर'' ही संज्ञा वापरतो. यामागे ऐतिहासिक कारणे आहेत.}
\subsection{$\lexists$ साठीचे नियम}
\begin{defish}
\AxiomC{$\Atom{\varphi}{t}$}
\RightLabel{\Intro{\lexists}}
\UnaryInfC{$\lexists[x][\Atom{\varphi}{x}]$}
\DisplayProof
\hfill
\AxiomC{$\lexists[x][\Atom{\varphi}{x}]$}
\AxiomC{[$\Atom{\varphi}{a}$]$^n$}
\DeduceC{$\chi$}
\DischargeRule{\Elim{\lexists}}{n}
\BinaryInfC{$\chi$}
\DisplayProof
\end{defish}
पुन्हा, $t$ हे बंद पद आहे आणि $a$ हा असा स्थिरांक आहे की तो
आधारविधान $\lexists[x][\varphi(x)]$ मध्ये, निष्कर्ष~$\chi$ मध्ये किंवा दोन
आधारविधानांनी समाप्त होणाऱ्या निष्पत्त्या मधील कोणत्याही
मुक्त न केलेली गृहीतकात ($\varphi(a)$ ही गृहीतके वगळता) येत नाही. आपण
$a$ ला \Elim{\lexists} अनुमानाचा \emph{आयगेन चर} म्हणतो.
\Intro{\lforall} आणि \Elim{\lexists} अनुमानांसाठी आयगेन चरावर लागू होणारे
निर्बंध वर स्वतंत्रपणे सांगितले आहेत. त्यांत संबंधित आधारविधान व निष्कर्ष,
तसेच आधारविधानांकडे नेणाऱ्या निष्पत्त्या मधील मुक्त न केलेली गृहीतके,
यांबाबत वेगवेगळे निर्बंध आहेत. या निर्बंधांना \emph{आयगेन चराची अट} म्हणतात.
\begin{explain}
ही प्रथा आठवा: $\varphi$ हे असे सूत्र असेल की त्यात
चल~$x$ मुक्त असेल, तर हे आपण~$\varphi(x)$ असे लिहून दर्शवतो.
त्याच संदर्भात, मग $\varphi(t)$ हा~$\Subst{\varphi}{t}{x}$ चा संक्षेप असतो.
त्यामुळे $\Intro\lexists$ नियम पुढीलप्रमाणेही लिहिता येईल:
\begin{prooftree}
\AxiomC{$\Subst{\varphi}{t}{x}$}
\RightLabel{\Intro{\lexists}}
\UnaryInfC{$\lexists[x][\varphi]$}
\end{prooftree}
लक्षात घ्या की~$\varphi$ मध्ये $t$ आधीपासून येऊ शकते; उदा., $\varphi$ हे
$\Atom{\Obj P}{t,x}$ असू शकते. म्हणून,~$\Atom{\Obj P}{t,t}$ पासून
$\lexists[x][\Atom{\Obj P}{t,x}]$ निष्पन्न करणे हे~$\Intro\lexists$
चे योग्य उपयोजन आहे---आधारविधानात~$t$ जेथे येते त्यांपैकी एक किंवा
अधिक ठिकाणी त्याऐवजी बद्ध चल~$x$ घालता येतो; सर्वच ठिकाणी
असा ``बदल'' करणे आवश्यक नाही. परंतु $\Intro\lforall$ आणि
$\Elim\lexists$ मधील आयगेन चराच्या अटीनुसार स्थिरांक~$a$ हा~$\varphi$
मध्ये येता कामा नये. म्हणून, $\Intro\lforall$ वापरून
$\Atom{\Obj P}{a,a}$ पासून $\lforall[x][\Atom{\Obj P}{a,x}]$ योग्य
रीतीने निष्पन्न करता येत नाही.
\end{explain}
\begin{explain}
\Intro{\lexists} आणि \Elim{\lforall} मध्ये पद~$t$ बंद
असण्याव्यतिरिक्त त्यावर आणखी कोणतेही निर्बंध नाहीत; म्हणून इतर कोणत्याही
अटीची काळजी करावी लागत नाही. याउलट, \Elim{\lexists} आणि
\Intro{\lforall} नियमांमध्ये आयगेन चराच्या अटीनुसार स्थिरांक~$a$ हा
निष्कर्षात किंवा कोणत्याही मुक्त न केलेली गृहीतकात कोठेही येता कामा नये.
ही पद्धत निर्दोष आहे, म्हणजे ती केवळ अशा मुक्त न केलेली गृहीतकांपासूनच
वाक्ये निष्पन्न करते की ज्यांच्यापासून ती तार्किकरीत्या निष्पन्न
होतात, याची खात्री करण्यासाठी ही अट आवश्यक आहे. ही अट नसती, तर पुढील
निष्पत्तीला परवानगी मिळाली असती:
\begin{prooftree}
  \AxiomC{$\lexists[x][\varphi(x)]$}
  \AxiomC{$\Discharge{\varphi(a)}{1}$}
  \RightLabel{*\Intro{\lforall}}
  \UnaryInfC{$\lforall[x][\varphi(x)]$}
  \RightLabel{\Elim{\lexists}}
  \BinaryInfC{$\lforall[x][\varphi(x)]$}
\end{prooftree}
परंतु, $\lexists[x][\varphi(x)] \Entails/ \lforall[x][\varphi(x)]$.
संख्यापकांच्या विलोपन नियमांमध्ये चलांच्या जागी केवळ बंद पदे
ठेवण्याची परवानगी असल्यामुळे, वाक्यांच्या संचापासून निष्पन्न करता
येणारे कोणतेही सूत्र हे स्वतः वाक्य असते, हे निष्पन्न होते.
\end{explain}
\section{निष्पत्ती}\label{fol:ntd:der:sec}
\begin{explain}
गृहीतक म्हणजे काय हे आपण सांगितले आहे आणि निगमन नियम दिले आहेत.
नैसर्गिक निगमनातील निष्पत्त्या यांपासून आगमनाने निर्माण होतात:
प्रत्येक निष्पत्ती ही एकतर स्वतंत्र गृहीतक असते किंवा एक, दोन वा तीन
निष्पत्त्या नंतर केलेले योग्य अनुमान असते.
\end{explain}
\begin{defn}[निष्पत्ती]
\emph{निष्पत्ती} म्हणजे गृहीतके~$\Gamma$
यांपासून वाक्य~$\varphi$ ची पुढील अटी पूर्ण करणारी वाक्यांची
सांत वृक्षरचना होय:
\begin{enumerate}
\item वृक्षातील सर्वांत वरची वाक्ये ही एकतर $\Gamma$ मध्ये असतात
  किंवा वृक्षातील एखाद्या अनुमानाने मुक्त केलेली असतात.
\item वृक्षातील सर्वांत खालचे वाक्य हे~$\varphi$ असते.
\item वृक्षाच्या तळाशी असलेले वाक्य~$\varphi$ वगळता, वृक्षातील प्रत्येक
  वाक्य हे एखाद्या निगमन नियमाच्या योग्य उपयोजनाचे आधारविधान असते
  आणि त्या उपयोजनाचा निष्कर्ष वृक्षात त्या वाक्याच्या थेट खाली येतो.
\end{enumerate}
मग $\varphi$ हा त्या निष्पत्तीचा \emph{निष्कर्ष} आणि $\Gamma$ ही तिची
मुक्त न केलेली गृहीतके आहेत, असे आपण म्हणतो.
जर~$\Gamma$ या गृहीतकांपासून $\varphi$ ची निष्पत्ती अस्तित्वात असेल, तर
$\varphi$ हे~$\Gamma$ पासून \emph{निष्पन्न करता येण्याजोगे} आहे असे म्हणतो; चिन्हांत:
$\Gamma \Proves \varphi$. प्रत्येक गृहीतक मुक्त केलेले असणारी $\varphi$
ची निष्पत्ती असेल, तर आपण~$\Proves \varphi$ असे लिहितो.
\end{defn}
\begin{ex}
प्रत्येक गृहीतक स्वतःच निष्पत्ती असते. म्हणून, उदा., एकटे $\varphi$ ही
निष्पत्ती आहे आणि एकटे $\psi$ सुद्धा. यांवर, समजा,
$\Intro{\land}$ नियम लागू करून नवी निष्पत्ती मिळवता येते:
\begin{prooftree}
\AxiomC{$\varphi$}
\AxiomC{$\psi$}
\RightLabel{\Intro{\land}}
\BinaryInfC{$\varphi \land \psi$}
\end{prooftree}
हे नियम सर्वसाधारण असतात: त्यांतील $\varphi$ आणि~$\psi$ यांच्या जागी कोणतीही
वाक्ये, उदा., $\chi$ आणि~$\theta$, ठेवता येतात. मग निष्कर्ष
$\chi \land \theta$ असेल, आणि म्हणून
\begin{prooftree}
\AxiomC{$\chi$}
\AxiomC{$\theta$}
\RightLabel{\Intro{\land}}
\BinaryInfC{$\chi \land \theta$}
\end{prooftree}
ही योग्य निष्पत्ती आहे. अर्थात, गृहीतकांची अदलाबदल करून $\theta$ ला
$\varphi$ ची आणि $\chi$ ला~$\psi$ ची भूमिका देता येते. त्यामुळे,
\begin{prooftree}
\AxiomC{$\theta$}
\AxiomC{$\chi$}
\RightLabel{\Intro{\land}}
\BinaryInfC{$\theta \land \chi$}
\end{prooftree}
हीसुद्धा योग्य निष्पत्ती आहे.
आता आपण दुसरा नियम, समजा, $\Intro{\lif}$ लागू करू शकतो. त्यामुळे सोपाधिक
विधानाचा निष्कर्ष काढता येतो आणि त्या विधानाच्या पूर्वांगाशी एकरूप असलेले
कोणतेही गृहीतक मुक्त करता येते. म्हणून पुढील दोन्या निष्पत्त्या योग्य आहेत:
\begin{prooftree}
\AxiomC{$\Discharge{\chi}{1}$}
\AxiomC{$\theta$}
\RightLabel{\Intro{\land}}
\BinaryInfC{$\chi \land \theta$}
\DischargeRule{\Intro{\lif}}{1}
\UnaryInfC{$\chi \lif (\chi \land \theta)$}
\DisplayProof\bottomAlignProof
\AxiomC{$\chi$}
\AxiomC{$\Discharge{\theta}{1}$}
\RightLabel{\Intro{\land}}
\BinaryInfC{$\chi \land \theta$}
\DischargeRule{\Intro{\lif}}{1}
\UnaryInfC{$\theta \lif (\chi \land \theta)$}
\end{prooftree}
त्या अनुक्रमे $\theta \Proves \chi \lif (\chi \land \theta)$ आणि
$\chi \Proves \theta \lif (\chi \land \theta)$ हे दाखवतात.
गृहीतके मुक्त करणे ही परवानगी आहे, आवश्यकता नाही, हे आठवा: गृहीतके मुक्त
करणे बंधनकारक नसते. विशेषतः, निष्पत्तीमध्ये ती गृहीतके नसली तरीही
नियम लागू करता येतो. उदाहरणार्थ, मुक्त करण्यासाठी~$\varphi$ हे गृहीतक
नसले तरी पुढील उपयोजन अनुज्ञात आहे:
\begin{prooftree}
  \AxiomC{$\psi$}
  \DischargeRule{\Intro{\lif}}{1}
  \UnaryInfC{$\varphi \lif \psi$}
\end{prooftree}
\end{ex}
\section{निष्पत्ती: उदाहरणे}\label{fol:ntd:pro:sec}
\begin{ex}
वाक्य $(\varphi \land \psi) \lif \varphi$ ची निष्पत्ती देऊ.
हवा असलेला निष्कर्ष निष्पत्तीच्या तळाशी लिहून आपण सुरुवात करतो.
\begin{prooftree}
\AxiomC{}
\UnaryInfC{$(\varphi\land \psi) \lif \varphi$}
\end{prooftree}
पुढे, या रूपाचे वाक्य कोणत्या प्रकारच्या अनुमानामुळे मिळू शकते हे
शोधावे लागते. निष्कर्षाचा मुख्य संकारक हा $\lif$ आहे; म्हणून
\Intro{\lif} नियम वापरून निष्कर्षापर्यंत पोहोचण्याचा प्रयत्न करू. काम पुढे
जाताना संबंधित गृहीतके लिहून ठेवणे आणि निगमन नियमांना लेबले देणे उत्तम;
त्यामुळे सिद्धतेच्या शेवटी सर्व गृहीतके मुक्त झाली आहेत का हे सहज
दिसते.
\begin{prooftree}
\AxiomC{$\Discharge{\varphi \land \psi}{1}$}
\DeduceC{$\varphi$}
\DischargeRule{\Intro{\lif}}{1}
\UnaryInfC{$(\varphi\land \psi) \lif \varphi$}
\end{prooftree}
आता गृहीतक $\varphi \land \psi$ पासून $\varphi$ पर्यंतच्या पायऱ्या भराव्या लागतील.
येथे हाताळण्यासाठी फक्त एकच संयोजक, $\land$, असल्यामुळे $\land$ चा विलोपन
नियम वापरलाच पाहिजे. त्यामुळे पुढील सिद्धता मिळते:
\begin{prooftree}
\AxiomC{$\Discharge{\varphi \land \psi}{1}$}
\RightLabel{\Elim{\land}}
\UnaryInfC{$\varphi$}
\DischargeRule{\Intro{\lif}}{1}
\UnaryInfC{$(\varphi\land \psi) \lif \varphi$}
\end{prooftree}
आता आपल्याकडे $(\varphi \land \psi) \lif \varphi$ ची योग्य निष्पत्ती आहे.
\end{ex}
\begin{ex}
आता $(\lnot \varphi \lor \psi) \lif (\varphi \lif \psi)$ ची निष्पत्ती देऊ.
हवा असलेला निष्कर्ष निष्पत्तीच्या तळाशी लिहून आपण सुरुवात करतो.
\begin{prooftree}
\AxiomC{}
\UnaryInfC{$(\lnot \varphi \lor \psi) \lif (\varphi \lif \psi)$}
\end{prooftree}
हा निष्कर्ष देऊ शकणारा तार्किक नियम शोधण्यासाठी निष्कर्षातील $\lnot$,
$\lor$ आणि $\lif$ हे तार्किक संयोजक पाहतो. सध्या $\lif$ ची पहिली आस्थितीच
महत्त्वाची आहे, कारण अंतिम क्रमवर्तीतील वाक्याचा तो मुख्य संकारक
आहे; $\lnot$, $\lor$ आणि $\lif$ ची दुसरी आस्थिती दुसऱ्या संयोजकाच्या
व्याप्तीत असल्यामुळे त्यांचा विचार नंतर करू. म्हणून आपण \Intro{\lif}
नियमाने सुरुवात करतो. त्याचे योग्य उपयोजन पुढीलप्रमाणे असले पाहिजे:
\begin{prooftree}
\AxiomC{$\Discharge{\lnot \varphi \lor \psi}{1}$}
\DeduceC{$\varphi \lif \psi$}
\DischargeRule{\Intro{\lif}}{1}
\UnaryInfC{$(\lnot \varphi \lor \psi) \lif (\varphi \lif \psi)$}
\end{prooftree}
पुढे जाण्यासाठी आता दोन पर्याय आहेत. एकतर तळापासून वरच्या दिशेने काम
चालू ठेवून \Intro{\lif} नियमाचे आणखी एक उपयोजन शोधता येईल, किंवा वरून
खाली काम करून \Elim{\lor} नियम लागू करता येईल. दुसरा पर्याय घेऊ.
\Elim{\lor} चे सर्वांत डावे आधारविधान म्हणून $\lnot \varphi \lor \psi$ हे
गृहीतक वापरू. \Elim{\lor} च्या योग्य उपयोजनासाठी इतर दोन्ही आधारविधाने
$\varphi \lif \psi$ या निष्कर्षाशी एकरूप असली पाहिजेत; मात्र प्रत्येकाला
अनुक्रमे आणखी एका गृहीतकापासून, म्हणजे $\lnot \varphi \lor \psi$ च्या दोन
विकल्पघटकांपैकी एकापासून, निष्पन्न करता येते. त्यामुळे आपली
निष्पत्ती पुढीलप्रमाणे दिसेल:
\begin{prooftree}
\AxiomC{$\Discharge{\lnot \varphi \lor \psi}{1}$}
\AxiomC{$\Discharge{\lnot \varphi}{2}$}
\DeduceC{$\varphi \lif \psi$}
\AxiomC{$\Discharge{\psi}{2}$}
\DeduceC{$\varphi \lif \psi$}
\DischargeRule{\Elim{\lor}}{2}
\TrinaryInfC{$\varphi \lif \psi$}
\DischargeRule{\Intro{\lif}}{1}
\UnaryInfC{$(\lnot \varphi \lor \psi) \lif (\varphi \lif \psi)$}
\end{prooftree}
उजवीकडील दोन्ही शाखांत आपल्याला $\varphi \lif \psi$ हे निष्पन्न करायचे आहे;
त्यासाठी \Intro{\lif} वापरणे सर्वांत योग्य आहे.
\begin{prooftree}
\AxiomC{$\Discharge{\lnot \varphi \lor \psi}{1}$}
\AxiomC{$\Discharge{\lnot \varphi}{2}, \Discharge{\varphi}{3}$}
\DeduceC{$\psi$}
\DischargeRule{\Intro{\lif}}{3}
\UnaryInfC{$\varphi \lif \psi$}
\AxiomC{$\Discharge{\psi}{2}, \Discharge{\varphi}{4}$}
\DeduceC{$\psi$}
\DischargeRule{\Intro{\lif}}{4}
\UnaryInfC{$\varphi \lif \psi$}
\DischargeRule{\Elim{\lor}}{2}
\TrinaryInfC{$\varphi \lif \psi$}
\DischargeRule{\Intro{\lif}}{1}
\UnaryInfC{$(\lnot \varphi \lor \psi) \lif (\varphi \lif \psi)$}
\end{prooftree}
निष्पत्तीच्या उरलेल्या दोन भागांसाठी, मधल्या भागात $\lnot \varphi$ आणि
$\varphi$ यांपासून, तर उजवीकडील भागात $\varphi$ आणि $\psi$ यांपासून, $\psi$ मिळवणाऱ्या
निष्पत्त्या हवी आहेत. यांपैकी पहिली आधी घेऊ. $\lnot \varphi$ आणि $\varphi$ ही
\Elim{\lnot} ची दोन आधारविधाने आहेत:
\begin{prooftree}
\AxiomC{$\Discharge{\lnot \varphi}{2}$}
\AxiomC{$\Discharge{\varphi}{3}$}
\RightLabel{\Elim{\lnot}}
\BinaryInfC{$\lfalse$}
\DeduceC{$\psi$}
\end{prooftree}
\FalseInt वापरून $\psi$ हा निष्कर्ष मिळवता येतो आणि शाखा पूर्ण करता येते.
\begin{prooftree}
\AxiomC{$\Discharge{\lnot \varphi \lor \psi}{1}$}
\AxiomC{$\Discharge{\lnot \varphi}{2}$}
\AxiomC{$\Discharge{\varphi}{3}$}
\RightLabel{\Intro{\lfalse}}
\BinaryInfC{$\lfalse$}
\RightLabel{\FalseInt}
\UnaryInfC{$\psi$}
\DischargeRule{\Intro{\lif}}{3}
\UnaryInfC{$\varphi \lif \psi$}
\AxiomC{$\Discharge{\psi}{2}, \Discharge{\varphi}{4}$}
\DeduceC{$\psi$}
\DischargeRule{\Intro{\lif}}{4}
\UnaryInfC{$\varphi \lif \psi$}
\DischargeRule{\Elim{\lor}}{2}
\TrinaryInfC{$\varphi \lif \psi$}
\DischargeRule{\Intro{\lif}}{1}
\UnaryInfC{$(\lnot \varphi \lor \psi) \lif (\varphi \lif \psi)$}
\end{prooftree}
आता सर्वांत उजवी शाखा पाहू. येथे हे लक्षात घेणे महत्त्वाचे आहे की
निष्पत्तीची व्याख्या
\emph{गृहीतके मुक्त करण्याची परवानगी देते}, पण
\emph{तसे करणे आवश्यक ठरवत नाही}. दुसऱ्या शब्दांत, $\varphi$ आणि $\psi$ या
गृहीतकांपैकी एकाचा उपयोग न करता दुसऱ्यापासून $\psi$ निष्पन्न करता येत असेल,
तर ते योग्य आहे. $\psi$ पासून $\psi$ हे निष्पन्न करणे तर क्षुल्लक आहे:
एकटे $\psi$ ही अशी निष्पत्ती असून कोणत्याही अनुमानाची गरज नाही. म्हणून
गृहीतक~$\varphi$ सरळ वगळता येते.
\begin{prooftree}
\AxiomC{$\Discharge{\lnot \varphi \lor \psi}{1}$}
\AxiomC{$\Discharge{\lnot \varphi}{2}$}
\AxiomC{$\Discharge{\varphi}{3}$}
\RightLabel{\Elim{\lnot}}
\BinaryInfC{$\lfalse$}
\RightLabel{\FalseInt}
\UnaryInfC{$\psi$}
\DischargeRule{\Intro{\lif}}{3}
\UnaryInfC{$\varphi \lif \psi$}
\AxiomC{$\Discharge{\psi}{2}$}
\RightLabel{\Intro{\lif}}
\UnaryInfC{$\varphi \lif \psi$}
\DischargeRule{\Elim{\lor}}{2}
\TrinaryInfC{$\varphi \lif \psi$}
\DischargeRule{\Intro{\lif}}{1}
\UnaryInfC{$(\lnot \varphi \lor \psi) \lif (\varphi \lif \psi)$}
\end{prooftree}
लक्षात घ्या की पूर्ण झालेल्या निष्पत्तीमध्ये सर्वांत उजवीकडील
\Intro{\lif} अनुमान प्रत्यक्षात कोणतेही गृहीतक मुक्त करत नाही.
\end{ex}
\begin{ex}
आतापर्यंत आपल्याला \FalseCl{} नियमाची गरज पडलेली नाही. हा नियम विशेष आहे,
कारण नियमाच्या निष्कर्षाचे उप-सूत्र नसलेले गृहीतक तो मुक्त करू देतो.
त्याचा \FalseInt{} नियमाशी जवळचा संबंध आहे. खरे तर \FalseInt{} नियम हे
\FalseCl{} नियमाचे विशेष प्रकरण आहे---``अंतःप्रज्ञावादी तर्कशास्त्र''
नावाचे एक तर्कशास्त्र आहे आणि त्यात केवळ \FalseInt{} अनुज्ञात आहे. इतर
कोणताही मार्ग चालत नसताना \FalseCl{} नियम हा अखेरचा उपाय असतो. उदाहरणार्थ,
आपल्याला $\varphi \lor \lnot \varphi$ हे निष्पन्न करायचे आहे असे समजा. नेहमीच्या
पद्धतीनुसार $\Intro{\lor}$ वापरून $\varphi \lor \lnot \varphi$ हे निष्पन्न
करण्याचा प्रयत्न करू. पण त्यासाठी कोणत्याही गृहीतकांशिवाय $\varphi$ किंवा
$\lnot \varphi$ यांपैकी एक निष्पन्न करावे लागेल, आणि ते शक्य नाही. अशा वेळी
\FalseCl{} उपयोगी पडतो!
\begin{prooftree}
  \AxiomC{$\Discharge{\lnot(\varphi \lor \lnot \varphi)}{1}$}
  \DeduceC{$\lfalse$}
  \DischargeRule{\FalseCl}{1}
  \UnaryInfC{$\varphi \lor \lnot \varphi$}
\end{prooftree}
आता $\lnot(\varphi \lor \lnot \varphi)$ पासून $\lfalse$ ची निष्पत्ती शोधायची
आहे. $\lfalse$ हा $\Elim{\lnot}$ चा निष्कर्ष असल्यामुळे तो नियम वापरून
पाहू:
\begin{prooftree}
  \AxiomC{$\Discharge{\lnot(\varphi \lor \lnot \varphi)}{1}$}
  \DeduceC{$\lnot \varphi$}
  \AxiomC{$\Discharge{\lnot(\varphi \lor \lnot \varphi)}{1}$}
  \DeduceC{$\varphi$}
  \RightLabel{\Elim{\lnot}}
  \BinaryInfC{$\lfalse$}
  \DischargeRule{\FalseCl}{1}
  \UnaryInfC{$\varphi \lor \lnot \varphi$}
\end{prooftree}
$\lnot \varphi$ ची निष्पत्ती शोधण्याच्या आपल्या पद्धतीनुसार
$\Intro{\lnot}$ चे उपयोजन करावे लागेल:
\begin{prooftree}
  \AxiomC{$\Discharge{\lnot(\varphi \lor \lnot \varphi)}{1}, \Discharge{\varphi}{2}$}
  \DeduceC{$\lfalse$}
  \DischargeRule{\Intro{\lnot}}{2}
  \UnaryInfC{$\lnot \varphi$}
  \AxiomC{$\Discharge{\lnot(\varphi \lor \lnot \varphi)}{1}$}
  \DeduceC{$\varphi$}
  \RightLabel{\Elim{\lnot}}
  \BinaryInfC{$\lfalse$}
  \DischargeRule{\FalseCl}{1}
  \UnaryInfC{$\varphi \lor \lnot \varphi$}
\end{prooftree}
येथे $\lnot(\varphi \lor \lnot \varphi)$ हे गृहीतक आणि आपल्या नव्या $\varphi$
गृहीतकापासून~$\Intro{\lor}$ ने निष्पन्न होणारे $\varphi \lor \lnot \varphi$
यांना $\Elim{\lnot}$ लावून $\lfalse$ सहज मिळवता येते:
\begin{prooftree}
  \AxiomC{$\Discharge{\lnot(\varphi \lor \lnot \varphi)}{1}$}
  \AxiomC{$\Discharge{\varphi}{2}$}
  \RightLabel{\Intro{\lor}}
  \UnaryInfC{$\varphi \lor \lnot \varphi$}
  \RightLabel{\Elim{\lnot}}
  \BinaryInfC{$\lfalse$}
  \DischargeRule{\Intro{\lnot}}{2}
  \UnaryInfC{$\lnot \varphi$}
  \AxiomC{$\Discharge{\lnot(\varphi \lor \lnot \varphi)}{1}$}
  \DeduceC{$\varphi$}
  \RightLabel{\Elim{\lnot}}
  \BinaryInfC{$\lfalse$}
  \DischargeRule{\FalseCl}{1}
  \UnaryInfC{$\varphi \lor \lnot \varphi$}
\end{prooftree}
उजव्या बाजूला हीच पद्धत वापरतो; फक्त तेथे $\varphi$ हे~\FalseCl ने मिळवतो:
\begin{prooftree}
  \AxiomC{$\Discharge{\lnot(\varphi \lor \lnot \varphi)}{1}$}
  \AxiomC{$\Discharge{\varphi}{2}$}
  \RightLabel{\Intro{\lor}}
  \UnaryInfC{$\varphi \lor \lnot \varphi$}
  \RightLabel{\Elim{\lnot}}
  \BinaryInfC{$\lfalse$}
  \DischargeRule{\Intro{\lnot}}{2}
  \UnaryInfC{$\lnot \varphi$}
  \AxiomC{$\Discharge{\lnot(\varphi \lor \lnot \varphi)}{1}$}
  \AxiomC{$\Discharge{\lnot \varphi}{3}$}
  \RightLabel{\Intro{\lor}}
  \UnaryInfC{$\varphi \lor \lnot \varphi$}
  \RightLabel{\Elim{\lnot}}
  \BinaryInfC{$\lfalse$}
  \DischargeRule{\FalseCl}{3}
  \UnaryInfC{$\varphi$}
  \RightLabel{\Elim{\lnot}}
  \BinaryInfC{$\lfalse$}
  \DischargeRule{\FalseCl}{1}
  \UnaryInfC{$\varphi \lor \lnot \varphi$}
\end{prooftree}
\end{ex}
\begin{prob}
पुढील गोष्टी दाखवणाऱ्या निष्पत्त्या द्या:
\begin{enumerate}
\item $\varphi \land (\psi \land \chi) \Proves (\varphi \land \psi) \land \chi$.
\item $\varphi \lor (\psi \lor \chi) \Proves (\varphi \lor \psi) \lor \chi$.
\item $\varphi \lif (\psi \lif \chi) \Proves \psi \lif (\varphi \lif \chi)$.
\item $\varphi \Proves \lnot\lnot \varphi$.
\end{enumerate}
\end{prob}
\begin{prob}
पुढील गोष्टी दाखवणाऱ्या निष्पत्त्या द्या:
\begin{enumerate}
\item $(\varphi \lor \psi) \lif \chi \Proves \varphi \lif \chi$.
\item $(\varphi \lif \chi) \land (\psi \lif \chi) \Proves (\varphi \lor \psi) \lif \chi$.
\item $\Proves \lnot(\varphi \land \lnot \varphi)$.
\item $\psi \lif \varphi \Proves \lnot \varphi \lif \lnot \psi$.
\item $\Proves (\varphi \lif \lnot \varphi) \lif \lnot \varphi$.
\item $\Proves \lnot(\varphi \lif \psi) \lif \lnot \psi$.
\item $\varphi \lif \chi \Proves \lnot (\varphi \land \lnot \chi)$.
\item $\varphi \land \lnot \chi \Proves \lnot (\varphi \lif \chi)$.
\item $\varphi \lor \psi, \lnot \psi \Proves \varphi$.
\item $\lnot \varphi \lor \lnot \psi \Proves \lnot(\varphi \land \psi)$.
\item $\Proves (\lnot \varphi \land \lnot \psi) \lif\lnot(\varphi \lor \psi)$.
\item $\Proves \lnot(\varphi \lor \psi) \lif (\lnot \varphi \land \lnot \psi)$.
\end{enumerate}
\end{prob}
\begin{prob}
पुढील गोष्टी दाखवणाऱ्या निष्पत्त्या द्या:
\begin{enumerate}
\item $\lnot(\varphi \lif \psi) \Proves \varphi$.
\item $\lnot(\varphi \land \psi) \Proves \lnot \varphi \lor \lnot \psi$.
\item $\varphi \lif \psi \Proves \lnot \varphi \lor \psi$.
\item $\Proves \lnot \lnot \varphi \lif \varphi$.
\item $\varphi \lif \psi, \lnot \varphi \lif \psi \Proves \psi$.
\item $(\varphi \land \psi) \lif \chi \Proves (\varphi \lif \chi) \lor (\psi \lif \chi)$.
\item $(\varphi \lif \psi) \lif \varphi \Proves \varphi$.
\item $\Proves (\varphi \lif \psi) \lor (\psi \lif \chi)$.
\end{enumerate}
(या सर्वांसाठी $\FalseCl$~नियम आवश्यक आहे.)
\end{prob}
\section{संख्यापकांसह निष्पत्ती}\label{fol:ntd:prq:sec}
\begin{ex}
संख्यापक हाताळताना आयगेन चराच्या अटीचा भंग होणार नाही याची खात्री करावी
लागते; त्यासाठी कधीकधी काही अनुमाने करण्याचा क्रम बदलून पाहावा लागतो.
सामान्यतः आयगेन चराच्या अटीच्या अधीन असलेले नियम आधी हाताळण्याचा प्रयत्न
उपयोगी ठरतो (पूर्ण सिद्धतेत ते अधिक खाली असतील).
सूत्र $\lexists[x][\lnot \varphi(x)] \lif \lnot \lforall[x][\varphi(x)]$
ची निष्पत्ती कशी द्यायची ते पाहू. नेहमीप्रमाणे सुरुवात करून लिहू:
\begin{prooftree}
\AxiomC{}
\UnaryInfC{$\lexists[x][\lnot \varphi(x)]\lif \lnot \lforall[x][\varphi(x)]$}
\end{prooftree}
ही शेवटची पायरी \Intro{\lif} नियमाने समर्थित होण्यासाठी काय आवश्यक आहे ते
प्रथम लिहू.
\begin{prooftree}
\AxiomC{$\Discharge{\lexists[x][\lnot \varphi(x)]}{1}$}
\DeduceC{$\lnot \lforall[x][\varphi(x)]$}
\DischargeRule{\Intro{\lif}}{1}
\UnaryInfC{$\lexists[x][\lnot \varphi(x)]\lif \lnot \lforall[x][\varphi(x)]$}
\end{prooftree}
$\lnot \lforall[x][\varphi(x)]$ वर लागू करता येईल असा कोणताही स्पष्ट नियम
नसल्यामुळे, \Elim{\lexists} नियम वापरता येईल अशा रीतीने निष्पत्ती
मांडू. येथे आयगेन चराच्या अटीकडे लक्ष देऊन, प्रमुख आधारविधान
$\lexists[x][\lnot \varphi(x)]$ किंवा निष्कर्ष यांत, तसेच त्यांच्याकडे नेणाऱ्या
निष्पत्तीतील इतर कोणत्याही मुक्त न केलेल्या गृहीतकात न येणारा अचर निवडावा
लागतो. (परंतु येथे कोणतेही स्थिरांक येत नसल्यामुळे कोणताही अचर चालेल.)
\begin{prooftree}
\AxiomC{$\Discharge{\lexists[x][\lnot \varphi(x)]}{1}$}
\AxiomC{$\Discharge{\lnot \varphi(a)}{2}$}
\DeduceC{$\lnot \lforall[x][\varphi(x)]$}
\DischargeRule{\Elim{\lexists}}{2}
\BinaryInfC{$\lnot \lforall[x][\varphi(x)]$}
\DischargeRule{\Intro{\lif}}{1}
\UnaryInfC{$\lexists[x][\lnot \varphi(x)] \lif \lnot \lforall[x][\varphi(x)]$}
\end{prooftree}
$\lnot \lforall[x][\varphi(x)]$ हे निष्पन्न करण्यासाठी \Intro{\lnot} नियम
वापरून पाहू: यासाठी, आवश्यक असल्यास $\lforall[x][\varphi(x)]$ हे अतिरिक्त गृहीतक वापरून,
व्याघात निष्पन्न करावा लागेल. अर्थात या व्याघातात $\lnot \varphi(a)$ हे गृहीतक
वापरले जाऊ शकते; ते \Elim{\lexists} अनुमानाने मुक्त केले जाईल. ही रचना
पुढीलप्रमाणे करता येते:
\begin{prooftree}
\AxiomC{$\Discharge{\lexists[x][\lnot \varphi(x)]}{1}$}
\AxiomC{$\Discharge{\lnot \varphi(a)}{2}, \Discharge{\lforall[x][\varphi(x)]}{3}$}
\DeduceC{$\lfalse$}
\DischargeRule{\Intro{\lnot}}{3}
\UnaryInfC{$\lnot \lforall[x][\varphi(x)]$}
\DischargeRule{\Elim{\lexists}}{2}
\BinaryInfC{$\lnot \lforall[x][\varphi(x)]$}
\DischargeRule{\Intro{\lif}}{1}
\UnaryInfC{$\lexists[x][\lnot \varphi(x)]\lif \lnot \lforall[x][\varphi(x)]$}
\end{prooftree}
व्याघात मिळण्याच्या आपण अगदी जवळ आलो आहोत. आयगेन चराची अट नसलेला
\Elim{\lforall} नियम येथे लागू करणे सर्वांत सोपे आहे. सार्वत्रिक
संख्यापकाने बद्ध असलेल्या~$x$ च्या जागी हवे ते कोणतेही बंद पद ठेवता येते;
व्याघातापर्यंत पोहोचण्यासाठी~$a$ चाच पुढे वापर करणे सर्वांत सयुक्तिक आहे.
\begin{prooftree}
\AxiomC{$\Discharge{\lexists[x][\lnot \varphi(x)]}{1}$}
\AxiomC{$\Discharge{\lnot \varphi(a)}{2}$}
\AxiomC{$\Discharge{\lforall[x][\varphi(x)]}{3}$}
\RightLabel{\Elim{\lforall}}
\UnaryInfC{$\varphi(a)$}
\RightLabel{\Elim{\lnot}}
\BinaryInfC{$\lfalse$}
\DischargeRule{\Intro{\lnot}}{3}
\UnaryInfC{$\lnot \lforall[x][\varphi(x)]$}
\DischargeRule{\Elim{\lexists}}{2}
\BinaryInfC{$\lnot \lforall[x][\varphi(x)]$}
\DischargeRule{\Intro{\lif}}{1}
\UnaryInfC{$\lexists[x][\lnot \varphi(x)]\lif \lnot \lforall[x][\varphi(x)]$}
\end{prooftree}
विशेषतः संख्यापक हाताळताना, आयगेन चराच्या अटीचा भंग झालेला नाही याची या
टप्प्यावर पुन्हा तपासणी करणे महत्त्वाचे आहे. आपण लावलेल्या नियमांपैकी फक्त
\Elim{\exists} हा नियम त्या अटीच्या अधीन होता. त्याचा आयगेन चर~$a$ प्रमुख
आधारविधानात किंवा निष्कर्षात येत नाही आणि नियमाने मुक्त केलेले तात्पुरते
गृहीतक वगळता संबंधित कोणत्याही मुक्त न केलेल्या गृहीतकातही येत नाही. म्हणून
ही योग्य निष्पत्ती आहे.
\end{ex}
\begin{ex}
कधीकधी इतर सूत्रांपासून सूत्र निष्पन्न करायचे असते. अशा वेळी
काही गृहीतके मुक्त न केलेली राहू शकतात. अंतिम ध्येयाबरोबरच गृहीतकांचीही
नोंद ठेवणे महत्त्वाचे आहे.
गृहीतके $\lexists[x][(\varphi(x) \land \psi(x))]$ आणि
$\lforall[x][(\psi(x) \lif \chi(x,b))]$ यांपासून सूत्र
$\lexists[x][\chi(x,b)]$ ची निष्पत्ती कशी द्यायची ते पाहू.
नेहमीप्रमाणे सुरुवात करून निष्कर्ष तळाशी लिहू.
\begin{prooftree}
\AxiomC{}
\UnaryInfC{$\lexists[x][\chi(x,b)]$}
\end{prooftree}
आपल्याकडे वापरण्यासाठी दोन आधारविधाने आहेत. पहिले वापरण्यासाठी, म्हणजे
$\lexists[x][(\varphi(x) \land \psi(x))]$ पासून
$\lexists[x][\chi(x, b)]$ ची निष्पत्ती शोधण्याचा प्रयत्न करण्यासाठी,
\Elim{\lexists} नियम वापरू. त्याला आयगेन चराची अट असल्यामुळे तो नियम आधी
लावू. मग पुढील मिळते:
\begin{prooftree}
\AxiomC{$\lexists[x][(\varphi(x) \land \psi(x))]$}
\AxiomC{$\Discharge{\varphi(a) \land \psi(a)}{1}$}
\DeduceC{$\lexists[x][\chi(x,b)]$}
\DischargeRule{\Elim{\lexists}}{1}
\BinaryInfC{$\lexists[x][\chi(x,b)]$}
\end{prooftree}
आपण वापरत असलेल्या दोन्ही गृहीतकांत~$\psi$ समान आहे. या टप्प्यावर
\Elim{\land} लावून~$\psi(a)$ वेगळे काढणे उपयोगी ठरेल.
\begin{prooftree}
\AxiomC{$\lexists[x][(\varphi(x) \land \psi(x)])$}
\AxiomC{$\Discharge{\varphi(a)
\land \psi(a)}{1}$}
\RightLabel{\Elim{\land}}
\UnaryInfC{$\psi(a)$}
\DeduceC{$\lexists[x][\chi(x,b)]$}
\DischargeRule{\Elim{\lexists}}{1}
\BinaryInfC{$\lexists[x][\chi(x,b)]$}
\end{prooftree}
वापरण्यासाठी असलेले दुसरे गृहीतक म्हणजे~$\lforall[x][(\psi(x) \lif
  \chi(x,b))]$. येथे आयगेन चराची अट नसल्यामुळे \Elim{\lforall} वापरून
स्थिरांक~$a$ हे $x$ च्या जागी ठेवता येते आणि
$\psi(a) \lif \chi(a, b)$ मिळते. आता आपल्याकडे $\psi(a) \lif \chi(a,b)$ आणि
$\psi(a)$ दोन्ही आहेत. यानंतर \Elim{\lif} नियमाचे सरळ उपयोजन करावे.
\begin{prooftree}
\AxiomC{$\lexists[x][(\varphi(x) \land \psi(x))]$}
\AxiomC{$\lforall[x][(\psi(x) \lif \chi(x,b))]$}
\RightLabel{\Elim{\lforall}}
\UnaryInfC{$\psi(a) \lif \chi(a,b)$}
\AxiomC{$\Discharge{\varphi(a)
\land \psi(a)}{1}$}
\RightLabel{\Elim{\land}}
\UnaryInfC{$\psi(a)$}
\RightLabel{\Elim{\lif}}
\BinaryInfC{$\chi(a,b)$}
\DeduceC{$\lexists[x][\chi(x,b)]$}
\DischargeRule{\Elim{\lexists}}{1}
\insertBetweenHyps{\hspace{-5em}}
\BinaryInfC{$\lexists[x][\chi(x,b)]$}
\end{prooftree}
आपण ध्येयाच्या अगदी जवळ आलो आहोत!{} \Intro{\lexists} चे एक उपयोजन केले
की ध्येय गाठले.
\begin{prooftree}
\AxiomC{$\lexists[x][(\varphi(x) \land \psi(x))]$}
\AxiomC{$\lforall[x][(\psi(x) \lif \chi(x,b))]$}
\RightLabel{\Elim{\lforall}}
\UnaryInfC{$\psi(a) \lif \chi(a,b)$}
\AxiomC{$\Discharge{\varphi(a)
\land \psi(a)}{1}$}
\RightLabel{\Elim{\land}}
\UnaryInfC{$\psi(a)$}
\RightLabel{\Elim{\lif}}
\BinaryInfC{$\chi(a,b)$}
\RightLabel{\Intro{\lexists}}
\UnaryInfC{$\lexists[x][\chi(x,b)]$}
\DischargeRule{\Elim{\lexists}}{1}
\insertBetweenHyps{\hspace{-5em}}
\BinaryInfC{$\lexists[x][\chi(x,b)]$}
\end{prooftree}
प्रत्येक पायरीवर आयगेन चराच्या अटींचा भंग झालेला नाही याची खात्री केल्यामुळे
ही योग्य निष्पत्ती आहे असा विश्वास ठेवता येतो.
\end{ex}
\begin{ex}
गृहीतके $\lforall[x][\varphi(x)] \lif \lexists[y][\psi(y)]$ आणि
$\lnot\lexists[y][\psi(y)]$ यांपासून सूत्र
$\lnot\lforall[x][\varphi(x)]$ ची निष्पत्ती द्या.
नेहमीप्रमाणे सुरुवात करून अपेक्षित सूत्र तळाशी लिहू.
\begin{prooftree}
\AxiomC{}
\UnaryInfC{$\lnot\lforall[x][\varphi(x)]$}
\end{prooftree}
निष्पत्तीची शेवटची ओळ नकरण आहे; म्हणून \Intro{\lnot} वापरून पाहू.
त्यासाठी व्याघात कसा निष्पन्न करायचा ते शोधावे लागेल.
\begin{prooftree}
\AxiomC{$\Discharge{\lforall[x][\varphi(x)]}{1}$}
\DeduceC{$\lfalse$}
\DischargeRule{\Intro{\lnot}}{1}
\UnaryInfC{$\lnot\lforall[x][\varphi(x)]$}
\end{prooftree}
येथपर्यंत सर्व ठीक आहे. \Elim{\lforall} वापरता येईल, पण त्याने ध्येयापर्यंत
पोहोचण्यास मदत होईल का हे स्पष्ट नाही. त्याऐवजी आपले एक गृहीतक वापरू.
$\lforall[x][\varphi(x)] \lif \lexists[y][\psi(y)]$ आणि
$\lforall[x][\varphi(x)]$ यांच्या आधारे \Elim{\lif} नियम वापरता येईल.
\begin{prooftree}
\AxiomC{$\lforall[x][\varphi(x)] \lif \lexists[y][\psi(y)]$}
\AxiomC{$\Discharge{\lforall[x][\varphi(x)]}{1}$}
\RightLabel{\Elim{\lif}}
\BinaryInfC{$\lexists[y][\psi(y)]$}
\DeduceC{$\lfalse$}
\DischargeRule{\Intro{\lnot}}{1}
\UnaryInfC{$\lnot\lforall[x][\varphi(x)]$}
\end{prooftree}
आता वापरण्यासाठी एक अखेरचे गृहीतक उरले आहे. \Elim{\lnot} वापरून
व्याघातापर्यंत पोहोचण्यास ते मदत करेल असे दिसते.
\begin{prooftree}
\AxiomC{$\lnot\lexists[y][\psi(y)]$}
\AxiomC{$\lforall[x][\varphi(x)] \lif \lexists[y][\psi(y)]$}
\AxiomC{$\Discharge{\lforall[x][\varphi(x)]}{1}$}
\RightLabel{\Elim{\lif}}
\BinaryInfC{$\lexists[y][\psi(y)]$}
\RightLabel{\Elim{\lnot}}
\BinaryInfC{$\lfalse$}
\DischargeRule{\Intro{\lnot}}{1}
\UnaryInfC{$\lnot\lforall[x][\varphi(x)]$}
\end{prooftree}
\end{ex}
\begin{prob}
पुढील गोष्टी दाखवणाऱ्या निष्पत्त्या द्या:
\begin{enumerate}
\item $\Proves (\lforall[x][\varphi(x)] \land \lforall[y][\psi(y)]) \lif
\lforall[z][(\varphi(z) \land \psi(z))]$.
\item $\Proves (\lexists[x][\varphi(x)] \lor \lexists[y][\psi(y)]) \lif
\lexists[z][(\varphi(z) \lor \psi(z))]$.
\item $\lforall[x][(\varphi(x) \lif \psi)] \Proves \lexists[y][\varphi(y)] \lif \psi$.
\item $\lforall[x][\lnot \varphi(x)] \Proves \lnot\lexists[x][\varphi(x)]$.
\item $\Proves \lnot\lexists[x][\varphi(x)] \lif \lforall[x][\lnot \varphi(x)]$.
\item $\Proves \lnot\lexists[x][\lforall[y][((\varphi(x,y) \lif \lnot
\varphi(y,y)) \land (\lnot \varphi(y,y) \lif \varphi(x,y)))]]$.
\end{enumerate}
\end{prob}
\begin{prob}
पुढील गोष्टी दाखवणाऱ्या निष्पत्त्या द्या:
\begin{enumerate}
\item $\Proves \lnot\lforall[x][\varphi(x)] \lif \lexists[x][\lnot\varphi(x)]$.
\item $(\lforall[x][\varphi(x)] \lif \psi) \Proves \lexists[y][(\varphi(y) \lif \psi)]$.
\item $\Proves \lexists[x][(\varphi(x) \lif \lforall[y][\varphi(y)])]$.
\end{enumerate}
(या सर्वांसाठी $\FalseCl$~नियम आवश्यक आहे.)
\end{prob}
\section{सिद्धता-उपपत्तीय संकल्पना}\label{fol:ntd:ptn:sec}
\begin{editorial}
या विभागात नैसर्गिक निगमनासाठी सिद्धतायोग्यता संबंध आणि सुसंगतता यांच्या
व्याख्या एकत्रित केल्या आहेत.
\end{editorial}
\begin{explain}
जशा आपण अनेक महत्त्वाच्या चिन्हार्थविषयक संकल्पना (वैधता, तार्किक निष्पन्नता,
पूर्ततायोग्यता) परिभाषित केल्या आहेत, तशाच आता त्यांना अनुरूप
\emph{सिद्धता-उपपत्तीय संकल्पना} परिभाषित करू. या संकल्पना संरचना मध्ये
वाक्यांची पूर्ति होते की नाही याच्या आधारे परिभाषित केलेल्या नसून,
काही वाक्यांची इतरांपासूनची निष्पन्नता किंवा
अनिष्पन्नता यांचा आधार घेऊन परिभाषित केल्या आहेत. या दोन प्रकारच्या
संकल्पना एकमेकांशी जुळतात, हा एक महत्त्वाचा शोध होता. त्या जुळतात, हाच
\emph{निर्दोषता प्रमेय} आणि \emph{संपूर्णता प्रमेय} यांचा आशय आहे.
\end{explain}
\begin{defn}[प्रमेये]
एखादे वाक्य~$\varphi$ हे \emph{प्रमेय} असते, जर नैसर्गिक निगमनात $\varphi$ ची
अशी निष्पत्ती असेल की तिच्यातील सर्व गृहीतके मुक्त आहेत.
$\varphi$ हे प्रमेय असेल तर आपण $\Proves \varphi$ लिहितो आणि ते प्रमेय नसेल तर
$\Proves/ \varphi$ लिहितो.
\end{defn}
\begin{defn}[निष्पन्नता]
वाक्य $\varphi$ हे वाक्यांच्या संच~$\Gamma$ \emph{पासून
निष्पन्न करता येण्याजोगे} असते, म्हणजे $\Gamma \Proves \varphi$, तेव्हा आणि केवळ तेव्हाच,
जेव्हा निष्कर्ष~$\varphi$ असणारी अशी निष्पत्ती असेल की तिच्यातील प्रत्येक
गृहीतक एकतर मुक्त आहे किंवा $\Gamma$ मध्ये आहे. $\varphi$ हे $\Gamma$
पासून निष्पन्न करता येण्याजोगे नसेल, तर आपण $\Gamma \Proves/ \varphi$ लिहितो.
\end{defn}
\begin{defn}[सुसंगतता]
वाक्यांचा संच~$\Gamma$ हा \emph{विसंगत} असतो तेव्हा आणि केवळ तेव्हाच,
जेव्हा $\Gamma \Proves \lfalse$. $\Gamma$ विसंगत नसेल, म्हणजे
$\Gamma \Proves/ \lfalse$ असेल, तर त्याला \emph{सुसंगत} म्हणतो.
\end{defn}
\begin{prop}[परावर्तिता]
\label{fol:ntd:ptn:prop:reflexivity}
$\varphi \in \Gamma$ असेल, तर $\Gamma \Proves \varphi$.
\end{prop}
\begin{proof}
केवळ $\varphi$ हे गृहीतक स्वतःच $\varphi$ ची अशी निष्पत्ती आहे की तिच्यातील
प्रत्येक मुक्त न केलेली गृहीतक (म्हणजे $\varphi$) $\Gamma$ मध्ये आहे.
\end{proof}
\begin{prop}[एकस्वनिकता]
\label{fol:ntd:ptn:prop:monotonicity}
$\Gamma \subseteq \Delta$ आणि $\Gamma \Proves \varphi$ असेल, तर $\Delta
\Proves \varphi$.
\end{prop}
\begin{proof}
$\Gamma$ पासून $\varphi$ ची कोणतीही निष्पत्ती ही $\Delta$ पासून $\varphi$ चीही
निष्पत्ती असते.
\end{proof}
\begin{prop}[संक्रमकता]
\label{fol:ntd:ptn:prop:transitivity}
$\Gamma \Proves \varphi$ आणि $\{\varphi\} \cup \Delta \Proves
\psi$ असेल, तर $\Gamma \cup \Delta \Proves \psi$.
\end{prop}
\begin{proof}
$\Gamma \Proves \varphi$ असेल, तर $\varphi$ ची अशी निष्पत्ती~$\delta_0$ असते
की तिच्यातील सर्व मुक्त न केलेली गृहीतके $\Gamma$ मध्ये असतात. $\{\varphi\}
\cup \Delta \Proves \psi$ असेल, तर $\psi$ ची अशी निष्पत्ती~$\delta_1$
असते की तिच्यातील सर्व मुक्त न केलेली गृहीतके $\{\varphi\} \cup \Delta$ मध्ये
असतात. आता पुढील रचना विचारात घ्या:
\begin{prooftree}
  \AxiomC{$\Delta, \Discharge{\varphi}{1}$}
  \RightLabel{$\delta_1$}
  \DeduceC{$\psi$}
  \DischargeRule{\Intro{\lif}}{1}
  \UnaryInfC{$\varphi \lif \psi$}
  \AxiomC{$\Gamma$}
  \RightLabel{$\delta_0$}
  \DeduceC{$\varphi$}
  \RightLabel{\Elim{\lif}}
  \BinaryInfC{$\psi$}
\end{prooftree}
आता सर्व मुक्त न केलेली गृहीतके $\Gamma \cup \Delta$ मध्ये आहेत; म्हणून
यावरून $\Gamma \cup \Delta \Proves \psi$ हे दिसते.
\end{proof}
$\Gamma = \{\varphi_1, \varphi_2, \ldots, \varphi_k\}$ हा सांत संच असेल, तेव्हा
$\Gamma \Proves \psi$ ऐवजी आपण $\varphi_1,\varphi_2,\ldots,\varphi_k \Proves \psi$ हे
सरलीकृत संकेतलेखन वापरू शकतो. विशेषतः, $\varphi \Proves \psi$ याचा अर्थ
$\{\varphi\} \Proves \psi$ असा होतो.
लक्षात घ्या की $\Gamma \Proves \varphi$ आणि $\varphi
\Proves \psi$ असेल, तर $\Gamma \Proves \psi$. तसेच $\varphi_1,
\dots, \varphi_n \Proves \psi$ असेल आणि प्रत्येक~$i$ साठी $\Gamma \Proves \varphi_i$
असेल, तर $\Gamma \Proves \psi$, हेसुद्धा यावरून मिळते.
\begin{prop}
\label{fol:ntd:ptn:prop:incons}
पुढील विधाने सममूल्य आहेत.
    \begin{enumerate}
      \item \( \Gamma \) विसंगत आहे.
      \item प्रत्येक वाक्य~\( {\varphi} \) साठी \( \Gamma \Proves {\varphi} \).
      \item एखाद्या वाक्य~\( {\varphi} \) साठी \( \Gamma \Proves {\varphi} \) आणि \( \Gamma \Proves \lnot {\varphi} \).
    \end{enumerate}
\end{prop}
\begin{proof}
सराव.
\end{proof}
\begin{prob}
\ref{fol:ntd:ptn:prop:incons} सिद्ध करा.
\end{prob}
\begin{prop}[संहतता]
  \label{fol:ntd:ptn:prop:proves-compact}
  \begin{enumerate}
  \item $\Gamma \Proves \varphi$ असेल, तर एखादा सांत उपसंच $\Gamma_0
    \subseteq \Gamma$ असा आहे की $\Gamma_0 \Proves \varphi$.
  \item $\Gamma$ चा प्रत्येक सांत उपसंच सुसंगत असेल, तर $\Gamma$ सुसंगत
    असतो.
  \end{enumerate}
\end{prop}
\begin{proof}
  \begin{enumerate}
    \item $\Gamma \Proves \varphi$ असेल, तर $\Gamma$ पासून $\varphi$ ची
      निष्पत्ती~$\delta$ असते. $\Gamma_0$ हा $\delta$ मधील
      मुक्त न केलेली गृहीतकांचा संच असू द्या. प्रत्येक निष्पत्ती सांत
      असल्यामुळे $\Gamma_0$ मधील वाक्यांची संख्या सांतच असते.
      म्हणून $\delta$ ही सांत~$\Gamma_0 \subseteq \Gamma$ पासून $\varphi$ ची
      निष्पत्ती आहे.
    \item (1) मध्ये $\varphi
      \ident \lfalse$ हे विशेष प्रकरण घेतल्यावर मिळणारे हे निषेधात्मक उलट
      विधान आहे.
  \end{enumerate}
\end{proof}
\section{निष्पन्नता आणि सुसंगतता}\label{fol:ntd:prv:sec}
आता आपण निष्पन्नता संबंधाचे अनेक गुणधर्म सिद्ध करू. ते स्वतंत्रपणेही
महत्त्वाचे आहेत, पण संपूर्णता प्रमेयाच्या सिद्धतेत प्रत्येकाची भूमिका असेल.
\begin{prop}\label{fol:ntd:prv:prop:provability-contr}
  जर $\Gamma \Proves \varphi$ आणि $\Gamma \cup \{\varphi\}$ विसंगत असेल, तर
  $\Gamma$ विसंगत आहे.
\end{prop}
\begin{proof}
$\Gamma$ पासून $\varphi$ ची निष्पत्ती~$\delta_1$ आणि $\Gamma \cup \{\varphi\}$
पासून $\lfalse$ ची निष्पत्ती~$\delta_2$ असू द्या. मग आपण पुढीलप्रमाणे
निष्पन्न करू शकतो:
\begin{prooftree}
\AxiomC{$\Gamma, \Discharge{\varphi}{1}$}
\RightLabel{$\delta_2$}
\DeduceC{$\lfalse$}
\DischargeRule{\Intro{\lnot}}{1}
\UnaryInfC{$\lnot \varphi$}
\AxiomC{$\Gamma$}
\RightLabel{$\delta_1$}
\DeduceC{$\varphi$}
\RightLabel{\Elim{\lnot}}
\BinaryInfC{$\lfalse$}
\end{prooftree}
नव्या निष्पत्तीमध्ये गृहीतक~$\varphi$ हे मुक्त आहे; म्हणून ही
$\Gamma$ पासूनची निष्पत्ती आहे.
\end{proof}
\begin{prop}
\label{fol:ntd:prv:prop:prov-incons}
$\Gamma \Proves \varphi$ तेव्हा आणि केवळ तेव्हाच, जेव्हा $\Gamma \cup \{\lnot \varphi\}$
विसंगत असते.
\end{prop}
\begin{proof}
प्रथम $\Gamma \Proves \varphi$ असे समजा; म्हणजे $\Gamma$ मधील
मुक्त न केलेली गृहीतकांपासून $\varphi$ ची निष्पत्ती~$\delta_0$ असते.
पुढीलप्रमाणे $\Gamma \cup \{\lnot \varphi\}$ पासून $\lfalse$ ची
निष्पत्ती मिळते:
\begin{prooftree}
  \AxiomC{$\lnot \varphi$}
  \AxiomC{$\Gamma$}
  \RightLabel{$\delta_0$}
  \DeduceC{$\varphi$}
  \RightLabel{\Elim{\lnot}}
  \BinaryInfC{$\lfalse$}
\end{prooftree}
आता $\Gamma \cup \{\lnot \varphi\}$ विसंगत आहे असे समजा आणि $\Gamma \cup
\{\lnot \varphi\}$ मधील मुक्त न केलेली गृहीतकांपासून $\lfalse$ ची संबंधित
निष्पत्ती~$\delta_1$ असू द्या. $\FalseCl$ वापरून पुढीलप्रमाणे केवळ
$\Gamma$ पासून $\varphi$ ची निष्पत्ती मिळते:
\begin{prooftree}
  \AxiomC{$\Gamma, \Discharge{\lnot \varphi}{1}$}
  \RightLabel{$\delta_1$}
  \DeduceC{$\lfalse$}
  \RightLabel{\FalseCl}
  \DischargeRule{\FalseCl}{1}
  \UnaryInfC{$\varphi$}
\end{prooftree}
\end{proof}
\begin{prob}
$\Gamma \Proves \lnot \varphi$ तेव्हा आणि केवळ तेव्हाच, जेव्हा $\Gamma \cup \{\varphi\}$
विसंगत असते, हे सिद्ध करा.
\end{prob}
\begin{prop}\label{fol:ntd:prv:prop:explicit-inc}
  जर $\Gamma \Proves \varphi$ आणि $\lnot \varphi \in \Gamma$ असेल, तर $\Gamma$
  विसंगत आहे.
\end{prop}
\begin{proof}
  समजा $\Gamma \Proves \varphi$ आणि $\lnot \varphi \in \Gamma$. मग $\Gamma$ पासून
  $\varphi$ ची निष्पत्ती~$\delta$ असते. $\Elim{\lnot}$ नियमाचा पुढील सोपा
  उपयोग विचारात घ्या:
  \begin{prooftree}
    \AxiomC{$\lnot \varphi$}
    \AxiomC{$\Gamma$}
    \RightLabel{$\delta$}
    \DeduceC{$\varphi$}
    \RightLabel{\Elim{\lnot}}
    \BinaryInfC{$\lfalse$}
  \end{prooftree}
  $\lnot \varphi \in \Gamma$ असल्यामुळे सर्व मुक्त न केलेली गृहीतके
  $\Gamma$ मध्ये आहेत. म्हणून यावरून $\Gamma \Proves \lfalse$ हे दिसते.
\end{proof}
\begin{prop}\label{fol:ntd:prv:prop:provability-exhaustive}
  $\Gamma \cup \{\varphi\}$ आणि $\Gamma \cup \{\lnot \varphi\}$ हे दोन्ही संच
  विसंगत असतील, तर $\Gamma$ विसंगत आहे.
\end{prop}
\begin{proof}
$\delta_1$ आणि $\delta_2$ या अनुक्रमे $\Gamma \cup \{ \varphi \}$ पासून
$\lfalse$ ची आणि $\Gamma \cup \{ \lnot \varphi \}$ पासून $\lfalse$ ची
निष्पत्त्या असू द्या. मग आपण पुढीलप्रमाणे निष्पन्न करू शकतो:
\begin{prooftree}
\AxiomC{$\Gamma, \Discharge{\lnot \varphi}{2}$}
\RightLabel{$\delta_2$}
\DeduceC{$\lfalse$}
\DischargeRule{\Intro{\lnot}}{2}
\UnaryInfC{$\lnot \lnot \varphi$}
\AxiomC{$\Gamma, \Discharge{\varphi}{1}$}
\RightLabel{$\delta_1$}
\DeduceC{$\lfalse$}
\DischargeRule{\Intro{\lnot}}{1}
\UnaryInfC{$\lnot \varphi$}
\RightLabel{\Elim{\lnot}}
\BinaryInfC{$\lfalse$}
\end{prooftree}
$\varphi$ आणि $\lnot \varphi$ ही गृहीतके मुक्त असल्यामुळे ही केवळ
$\Gamma$ पासून $\lfalse$ ची निष्पत्ती आहे. म्हणून $\Gamma$ विसंगत
आहे.
\end{proof}
\section{निष्पन्नता आणि विधानीय संयोजक}\label{fol:ntd:ppr:sec}
\begin{explain}
  नैसर्गिक निगमनाचा निष्पन्नता संबंध~$\Proves$ हा विधानीय संयोजकांशी
  संबंधित काही मूलभूत तथ्ये सिद्ध करण्यासाठी पुरेसा सामर्थ्यवान आहे; उदाहरणार्थ,
  $\varphi \land \psi
  \Proves \varphi$ आणि $\varphi, \varphi \lif \psi \Proves \psi$ (विध्यात्मक अनुमान). ही
  तथ्ये संपूर्णता प्रमेयाच्या सिद्धतेसाठी आवश्यक आहेत.
\end{explain}
\begin{prop}\label{fol:ntd:ppr:prop:provability-land}
  \begin{enumerate}
  \item \label{fol:ntd:ppr:prop:provability-land-left} $\varphi \land \psi \Proves
    \varphi$ आणि $\varphi \land \psi \Proves \psi$ ही दोन्ही तथ्ये सत्य आहेत.
  \item \label{fol:ntd:ppr:prop:provability-land-right} $\varphi, \psi \Proves \varphi \land \psi$.
  \end{enumerate}
\end{prop}
\begin{proof}
  \begin{enumerate}
  \item आपण पुढील दोन्ही निष्पन्न करू शकतो:
    \begin{prooftree}
      \AxiomC{$\varphi \land \psi$}
      \RightLabel{\Elim{\land}}
      \UnaryInfC{$\varphi$}
      \DisplayProof\qquad\bottomAlignProof
      \AxiomC{$\varphi \land \psi$}
      \RightLabel{\Elim{\land}}
      \UnaryInfC{$\psi$}
    \end{prooftree}
  \item आपण पुढील निष्पन्न करू शकतो:
    \begin{prooftree}
      \AxiomC{$\varphi$}
      \AxiomC{$\psi$}
      \RightLabel{\Intro{\land}}
      \BinaryInfC{$\varphi \land \psi$}
    \end{prooftree}
  \end{enumerate}
\end{proof}
\begin{prop}\label{fol:ntd:ppr:prop:provability-lor}
  \begin{enumerate}
  \item $\varphi \lor \psi, \lnot \varphi, \lnot \psi$ विसंगत आहे.
  \item $\varphi \Proves \varphi \lor \psi$ आणि $\psi \Proves \varphi \lor \psi$ ही दोन्ही
    तथ्ये सत्य आहेत.
  \end{enumerate}
\end{prop}
\begin{proof}
  \begin{enumerate}
  \item पुढील निष्पत्ती विचारात घ्या:
    \begin{prooftree}
      \AxiomC{$\varphi \lor \psi$}
      \AxiomC{$\lnot \varphi$}
      \AxiomC{$\Discharge{\varphi}{1}$}
      \RightLabel{\Elim{\lnot}}
      \BinaryInfC{$\lfalse$}
      \AxiomC{$\lnot \psi$}
      \AxiomC{$\Discharge{\psi}{1}$}
      \RightLabel{\Elim{\lnot}}
      \BinaryInfC{$\lfalse$}
      \DischargeRule{\Elim{\lor}}{1}
      \TrinaryInfC{$\lfalse$}
    \end{prooftree}
    ही $\varphi \lor \psi$, $\lnot \varphi$ आणि $\lnot \psi$ या मुक्त न केलेली
    गृहीतकांपासून $\lfalse$ ची निष्पत्ती आहे.
  \item आपण पुढील दोन्ही निष्पन्न करू शकतो:
    \begin{prooftree}
      \AxiomC{$\varphi$}
      \RightLabel{\Intro{\lor}}
      \UnaryInfC{$\varphi \lor \psi$}
      \DisplayProof\qquad\bottomAlignProof
      \AxiomC{$\psi$}
      \RightLabel{\Intro{\lor}}
      \UnaryInfC{$\varphi \lor \psi$}
    \end{prooftree}
  \end{enumerate}
\end{proof}
\begin{prop}\label{fol:ntd:ppr:prop:provability-lif}
  \begin{enumerate}
  \item \label{fol:ntd:ppr:prop:provability-lif-left} $\varphi, \varphi \lif \psi \Proves \psi$.
  \item \label{fol:ntd:ppr:prop:provability-lif-right}
    $\lnot \varphi \Proves \varphi \lif \psi$ आणि $\psi \Proves \varphi \lif \psi$ ही दोन्ही
    तथ्ये सत्य आहेत.
  \end{enumerate}
\end{prop}
\begin{proof}
  \begin{enumerate}
  \item आपण पुढील निष्पन्न करू शकतो:
    \begin{prooftree}
      \AxiomC{$\varphi \lif \psi$}
      \AxiomC{$\varphi$}
      \RightLabel{\Elim{\lif}}
      \BinaryInfC{$\psi$}
    \end{prooftree}
  \item हे पुढील दोन निष्पत्त्या दाखवतात:
    \begin{prooftree}
      \AxiomC{$\lnot \varphi$}
      \AxiomC{$\Discharge{\varphi}{1}$}
      \RightLabel{\Elim{\lnot}}
      \BinaryInfC{$\lfalse$}
      \RightLabel{\FalseInt}
      \UnaryInfC{$\psi$}
      \DischargeRule{\Intro{\lif}}{1}
      \UnaryInfC{$\varphi \lif \psi$}
      \DisplayProof\qquad\bottomAlignProof
      \AxiomC{$\psi$}
      \RightLabel{\Intro{\lif}}
      \UnaryInfC{$\varphi \lif \psi$}
    \end{prooftree}
    लक्षात घ्या की $\Intro{\lif}$ गृहीतक~$\varphi$ हे मुक्त करू शकतो;
    पण ते करणे आवश्यक नाही.
  \end{enumerate}
\end{proof}
\section{निष्पन्नता आणि संख्यापक}\label{fol:ntd:qpr:sec}
\begin{explain}
  संपूर्णता प्रमेयासाठी या विभागात स्थापित केलेली~$\Proves$ संबंधी तथ्ये
  नैसर्गिक निगमनाच्या नियमांनी निष्पन्न होतात, हेही आवश्यक आहे.
\end{explain}
\begin{thm}
\label{fol:ntd:qpr:thm:strong-generalization} जर $c$ हे $\Gamma$ किंवा $\varphi(x)$ मध्ये
न येणारे स्थिर असेल आणि $\Gamma \Proves \varphi(c)$, तर $\Gamma
\Proves \lforall[x][\varphi(x)]$.
\end{thm}
\begin{proof}
$\delta$ ही $\Gamma$ पासून $\varphi(c)$ ची निष्पत्ती असू द्या.
\Intro{\lforall} अनुमान जोडल्यावर आपल्याला $\lforall[x][\varphi(x)]$ ची
निष्पत्ती मिळते. $c$ हे $\Gamma$ किंवा $\varphi(x)$ मध्ये येत नसल्यामुळे
आयगेन चराची अट पूर्ण होते.
\end{proof}
\begin{prop}
\label{fol:ntd:qpr:prop:provability-quantifiers}
\begin{enumerate}
\item $\varphi(t) \Proves \lexists[x][\varphi(x)]$.
\item $\lforall[x][\varphi(x)] \Proves \varphi(t)$.
\end{enumerate}
\end{prop}
\begin{proof}
\begin{enumerate}
\item पुढील ही~$\lexists[x][\varphi(x)]$ ची $\varphi(t)$ पासूनची
  निष्पत्ती आहे:
\begin{prooftree}
\AxiomC{$\varphi(t)$}
\RightLabel{\Intro{\lexists}}
\UnaryInfC{$\lexists[x][\varphi(x)]$}
\end{prooftree}
\item पुढील ही~$\varphi(t)$ ची $\lforall[x][\varphi(x)]$ पासूनची
  निष्पत्ती आहे:
\begin{prooftree}
\AxiomC{$\lforall[x][\varphi(x)]$}
\RightLabel{\Elim{\lforall}}
\UnaryInfC{$\varphi(t)$}
\end{prooftree}
\end{enumerate}
\end{proof}
\section{निर्दोषता}\label{fol:ntd:sou:sec}
\begin{explain}
नैसर्गिक निगमनासारखी निष्पत्ती-पद्धत \emph{निर्दोष} असते, जर ती
प्रत्यक्षात निष्पन्न न होणाऱ्या गोष्टी निष्पन्न करू शकत नसेल. त्यामुळे
निर्दोषता हा निष्पत्ती-पद्धतींसाठी हमी दिलेल्या सुरक्षिततेचा एक प्रकार
आहे. कोणता सिद्धता-उपपत्तीय गुणधर्म विचारात आहे यावर अवलंबून, उदाहरणार्थ,
आपल्याला पुढील गोष्टी हव्या असतात:
\begin{enumerate}
\item प्रत्येक निष्पन्न करता येण्याजोगे वाक्य हे वैध असते;
\item एखादे वाक्य इतर काही वाक्यांपासून निष्पन्न करता येण्याजोगे असेल, तर ते
  त्यांचा तार्किक परिणामही असते;
\item वाक्यांचा संच विसंगत असेल, तर तो अपूर्ततायोग्य असतो.
\end{enumerate}
हे निष्पत्ती-पद्धतीचे महत्त्वाचे गुणधर्म आहेत. यांपैकी एखादा गुणधर्म
खरा नसेल, तर निष्पत्ती-पद्धतीत दोष आहे---ती खूप जास्त गोष्टी
निष्पन्न करेल. म्हणून निष्पत्ती-पद्धतीची निर्दोषता सिद्ध करणे अत्यंत
महत्त्वाचे आहे.
\end{explain}
\begin{thm}[निर्दोषता]
\label{fol:ntd:sou:thm:soundness}
$\varphi$ हे मुक्त न केलेली गृहीतके~$\Gamma$ यांच्यापासून निष्पन्न करता येण्याजोगे असेल, तर
$\Gamma \Entails \varphi$.
\end{thm}
\begin{proof}
$\delta$ ही $\varphi$ ची निष्पत्ती असू द्या. $\delta$ मधील अनुमानांच्या
संख्येवर विगमनाने पुढे जाऊ.
विगमनाच्या आधारपायरीसाठी अनुमानांची संख्या~$0$ असेल तेव्हा दावा सिद्ध करू.
या प्रकरणात $\delta$ मध्ये फक्त एकच वाक्य~$\varphi$, म्हणजे एक गृहीतक,
असते. ते गृहीतक मुक्त न केलेली असते, कारण गृहीतके फक्त अनुमानांद्वारे
मुक्त होऊ शकतात आणि येथे कोणतेही अनुमान नाही. म्हणून सिद्धतेतील
सर्व मुक्त न केलेली गृहीतके पूर्ण करणारी कोणतीही
संरचना~$\Struct{M}$
$\varphi$ हेदेखील पूर्ण करते.
आता विगमन पायरी पाहू. $\delta$ मध्ये $n$ अनुमाने आहेत असे समजा. सर्वांत
खालच्या अनुमानाची आधारविधाने अशा उप-निष्पत्त्या वापरून निष्पन्न केली
आहेत की प्रत्येकामध्ये $n$ पेक्षा कमी अनुमाने आहेत. विगमन गृहीतक असे धरू:
सर्वांत खालच्या अनुमानाची आधारविधाने, त्या आधारविधानांवर संपणाऱ्या
उप-निष्पत्त्यांच्या मुक्त न केलेली गृहीतकांपासून निष्पन्न होतात. संपूर्ण
सिद्धतेच्या मुक्त न केलेली गृहीतकांपासून निष्कर्ष~$\varphi$ निष्पन्न होतो, हे
आपल्याला दाखवायचे आहे.
सर्वांत खालच्या अनुमानाच्या प्रकारानुसार प्रकरणे वेगळी करू. प्रथम एकच
आधारविधान असलेली संभाव्य अनुमाने पाहू.
\begin{enumerate}
\item शेवटचे अनुमान \Intro{\lnot} आहे असे समजा: निष्पत्तीचे रूप
  पुढीलप्रमाणे आहे:
  \begin{prooftree}
    \AxiomC{$\Gamma, \Discharge{\varphi}{n}$}
    \RightLabel{$\delta_1$}
    \DeduceC{$\lfalse$}
    \DischargeRule{\Intro{\lnot}}{n}
    \UnaryInfC{$\lnot \varphi$}
  \end{prooftree}
  विगमन गृहीतकानुसार $\lfalse$ हे $\delta_1$ च्या मुक्त न केलेली
  गृहीतकांपासून~$\Gamma \cup \{\varphi\}$ निष्पन्न होते. एखादी
  संरचना~$\Struct{M}$
  विचारात घ्या. जर $\Sat{M}{\Gamma}$, तर
  $\Sat{M}{\lnot \varphi}$, हे दाखवायचे आहे.
  विसंगतीद्वारे सिद्धतेसाठी असे समजा की
  $\Sat{M}{\Gamma}$, पण
  $\Sat/{M}{\lnot \varphi}$, म्हणजे
  $\Sat{M}{\varphi}$. याचा अर्थ
  $\Sat{M}{\Gamma \cup \{\varphi\}}$ असा होईल. हे आपल्या
  विगमन गृहीतकाच्या विरुद्ध आहे. म्हणून
  $\Sat{M}{\lnot \varphi}$.
\item शेवटचे अनुमान \Elim{\land} आहे. याचे दोन प्रकार आहेत: $\varphi \land
  \psi$ या आधारविधानापासून $\varphi$ किंवा $\psi$ निष्पन्न केले जाऊ शकते. पहिले
  प्रकरण पाहू. निष्पत्ती~$\delta$ पुढीलप्रमाणे दिसते:
  \begin{prooftree}
    \AxiomC{$\Gamma$}
    \RightLabel{$\delta_1$}
    \DeduceC{$\varphi \land \psi$}
    \RightLabel{\Elim{\land}}
    \UnaryInfC{$\varphi$}
  \end{prooftree}
  विगमन गृहीतकानुसार $\varphi \land \psi$ हे $\delta_1$ च्या
  मुक्त न केलेली गृहीतकांपासून~$\Gamma$ निष्पन्न होते. एखादी
  संरचना~$\Struct{M}$ विचारात घ्या.
  जर $\Sat{M}{\Gamma}$, तर
  $\Sat{M}{\varphi}$, हे दाखवायचे आहे.
  $\Sat{M}{\Gamma}$ असे समजा. आपल्या विगमन
  गृहीतकानुसार ($\Gamma \Entails \varphi \land \psi$),
  $\Sat{M}{\varphi \land \psi}$ हे आपल्याला माहीत आहे.
  व्याख्येनुसार, $\Sat{M}{\varphi \land \psi}$ तेव्हा आणि
  केवळ तेव्हाच, जेव्हा $\Sat{M}{\varphi}$ आणि
  $\Sat{M}{\psi}$. ($\varphi \land \psi$ पासून $\psi$
  निष्पन्न करण्याचे प्रकरणही असेच हाताळता येते.)
\item शेवटचे अनुमान \Intro{\lor} आहे. याचे दोन प्रकार आहेत: $\varphi$ या
  आधारविधानापासून किंवा $\psi$ या आधारविधानापासून $\varphi \lor \psi$ निष्पन्न
  केले जाऊ शकते. पहिले प्रकरण पाहू. निष्पत्तीचे रूप पुढीलप्रमाणे आहे:
  \begin{prooftree}
    \AxiomC{$\Gamma$}
    \RightLabel{$\delta_1$}
    \DeduceC{$\varphi$}
    \RightLabel{\Intro{\lor}}
    \UnaryInfC{$\varphi \lor \psi$}
  \end{prooftree}
  विगमन गृहीतकानुसार $\varphi$ हे $\delta_1$ च्या मुक्त न केलेली
  गृहीतकांपासून~$\Gamma$ निष्पन्न होते. एखादी
  संरचना~$\Struct{M}$
  विचारात घ्या. जर $\Sat{M}{\Gamma}$, तर
  $\Sat{M}{\varphi \lor \psi}$, हे दाखवायचे आहे.
  $\Sat{M}{\Gamma}$ असे समजा; मग
  $\Gamma \Entails \varphi$ (विगमन गृहीतक) असल्यामुळे
  $\Sat{M}{\varphi}$. म्हणून
  $\Sat{M}{\varphi \lor \psi}$ हेदेखील असलेच पाहिजे.
  ($\psi$ पासून $\varphi \lor \psi$ निष्पन्न करण्याचे प्रकरणही असेच हाताळता येते.)
\item शेवटचे अनुमान \Intro{\lif} आहे. $\varphi$ हे गृहीतक आणि $\psi$ हा निष्कर्ष
  असलेल्या उपसिद्धतेपासून $\varphi \lif \psi$ निष्पन्न केले जाते, म्हणजे:
  \begin{prooftree}
    \AxiomC{$\Gamma, \Discharge{\varphi}{n}$}
    \RightLabel{$\delta_1$}
    \DeduceC{$\psi$}
    \DischargeRule{\Intro{\lif}}{n}
    \UnaryInfC{$\varphi \lif \psi$}
  \end{prooftree}
  विगमन गृहीतकानुसार $\psi$ हे $\delta_1$ च्या मुक्त न केलेली
  गृहीतकांपासून निष्पन्न होते, म्हणजे $\Gamma \cup \{\varphi\} \Entails \psi$.
  एखादी
  संरचना~$\Struct{M}$
  विचारात घ्या. शेवटच्या अनुमानात $\varphi$ चे विसर्जन झाल्यामुळे $\delta$
  ची मुक्त न केलेली गृहीतके फक्त $\Gamma$ आहेत. म्हणून
  $\Gamma \Entails \varphi \lif \psi$ हे दाखवायचे आहे. विसंगतीद्वारे सिद्धतेसाठी
  असे समजा की एखाद्या
  संरचना~$\Struct{M}$
  साठी $\Sat{M}{\Gamma}$, पण
  $\Sat/{M}{\varphi \lif \psi}$. मग
  $\Sat{M}{\varphi}$ आणि
  $\Sat/{M}{\psi}$. पण गृहीतकानुसार $\psi$ हा
  $\Gamma \cup \{\varphi\}$ चा तार्किक परिणाम आहे, म्हणजे
  $\Sat{M}{\psi}$; ही विसंगती आहे. म्हणून
  $\Gamma \Entails \varphi \lif \psi$.
\item शेवटचे अनुमान \FalseInt आहे. येथे $\delta$ पुढीलप्रमाणे संपते:
  \begin{prooftree}
    \AxiomC{$\Gamma$}
    \RightLabel{$\delta_1$}
    \DeduceC{$\lfalse$}
    \RightLabel{\FalseInt}
    \UnaryInfC{$\varphi$}
  \end{prooftree}
  विगमन गृहीतकानुसार $\Gamma \Entails \lfalse$. आपल्याला
  $\Gamma \Entails \varphi$ हे दाखवायचे आहे. तसे नसल्याचे समजा; मग एखाद्या
  $\Struct{M}$ साठी
    $\Sat{M}{\Gamma}$ आणि
    $\Sat/{M}{\varphi}$. पण नेहमीच
    $\Sat/{M}{\lfalse}$ असल्यामुळे याचा अर्थ
    $\Gamma \Entails/ \lfalse$ असा होईल, जो विगमन गृहीतकाच्या विरुद्ध आहे.
\item शेवटचे अनुमान \FalseCl आहे: सराव.
\item शेवटचे अनुमान \Intro{\lforall} आहे. मग $\delta$ चे रूप पुढीलप्रमाणे आहे:
  \begin{prooftree}
    \AxiomC{$\Gamma$}
    \RightLabel{$\delta_1$}
    \DeduceC{$\varphi(a)$}
    \RightLabel{\Intro{\lforall}}
    \UnaryInfC{$\lforall[x][\varphi(x)]$}
  \end{prooftree}
  विगमन गृहीतकानुसार आधारविधान~$\varphi(a)$ हे मुक्त न केलेली
  गृहीतके~$\Gamma$ यांचा तार्किक परिणाम आहे.
  $\Sat{M}{\Gamma}$ अशी एखादी रचना~$\Struct{M}$ विचारात घ्या. आपल्याला
  $\Sat{M}{\lforall[x][\varphi(x)]}$ हे दाखवायचे आहे.
  $\lforall[x][\varphi(x)]$ हे वाक्य असल्यामुळे याचा अर्थ प्रत्येक
  चल-विन्यास~$s$ साठी $\Sat{M}{\varphi(x)}[s]$ दाखवायचे आहे
  (\ref{fol:syn:ass:prop:sat-quant}). $\Gamma$ मध्ये केवळ वाक्ये असल्यामुळे
  \ref{fol:syn:sat:defn:satisfaction} नुसार प्रत्येक $\psi \in \Gamma$ साठी
  $\Sat{M}{\psi}[s]$. $\Struct{M'}$ ही $\Struct{M}$ सारखीच रचना असू द्या,
  फक्त $\Assign{a}{M'} = s(x)$. $a$ हे $\Gamma$ मध्ये येत नसल्यामुळे
  \ref{fol:syn:ext:cor:extensionality-sent} नुसार
  $\Sat{M'}{\Gamma}$. $\Gamma \Entails \varphi(a)$ असल्यामुळे
  $\Sat{M'}{\varphi(a)}$. $\varphi(a)$ हे वाक्य असल्यामुळे
  \ref{fol:syn:ass:prop:sentence-sat-true} नुसार
  $\Sat{M'}{\varphi(a)}[s]$. \ref{fol:syn:ext:prop:ext-formulas} नुसार
  $\Sat{M'}{\varphi(x)}[s]$ तेव्हा आणि केवळ तेव्हाच, जेव्हा
  $\Sat{M'}{\varphi(a)}$ (लक्षात घ्या की $\varphi(a)$ म्हणजे
  $\Subst{\varphi(x)}{a}{x}$). म्हणून $\Sat{M'}{\varphi(x)}[s]$. $a$ हे $\varphi(x)$
  मध्ये येत नसल्यामुळे \ref{fol:syn:ext:prop:extensionality} नुसार
  $\Sat{M}{\varphi(x)}[s]$. पण $s$ हा स्वेच्छ चल-विन्यास होता, म्हणून
  $\Sat{M}{\lforall[x][\varphi(x)]}$.
\item शेवटचे अनुमान \Intro{\lexists} आहे: सराव.
\item शेवटचे अनुमान \Elim{\forall} आहे: सराव.
\end{enumerate}
आता अनेक आधारविधाने असलेली संभाव्य अनुमाने पाहू:
\Elim{\lor}, \Intro{\land}, \Elim{\lif} आणि
  \Elim{\lexists}.
\begin{enumerate}
\item शेवटचे अनुमान \Intro{\land} आहे. $\varphi$ आणि $\psi$ या आधारविधानांपासून
  $\varphi \land \psi$ निष्पन्न केले जाते आणि $\delta$ चे रूप पुढीलप्रमाणे आहे:
  \begin{prooftree}
    \AxiomC{$\Gamma_1$}
    \RightLabel{$\delta_1$}
    \DeduceC{$\varphi$}
    \AxiomC{$\Gamma_2$}
    \RightLabel{$\delta_2$}
    \DeduceC{$\psi$}
    \RightLabel{\Intro{\land}}
    \BinaryInfC{$\varphi \land \psi$}
  \end{prooftree}
  विगमन गृहीतकानुसार $\varphi$ हे $\delta_1$ च्या मुक्त न केलेली
  गृहीतकांपासून~$\Gamma_1$ निष्पन्न होते आणि $\psi$ हे $\delta_2$ च्या
  मुक्त न केलेली गृहीतकांपासून~$\Gamma_2$ निष्पन्न होते. $\delta$ ची
  मुक्त न केलेली गृहीतके $\Gamma_1 \cup \Gamma_2$ आहेत; म्हणून
  $\Gamma_1 \cup \Gamma_2 \Entails \varphi \land \psi$ हे दाखवायचे आहे.
  $\Sat{M}{\Gamma_1 \cup \Gamma_2}$ असलेली एखादी
  संरचना~$\Struct{M}$
  विचारात घ्या. $\Sat{M}{\Gamma_1}$ आणि
  $\Gamma_1 \Entails \varphi$ असल्यामुळे
  $\Sat{M}{\varphi}$ असलेच पाहिजे; तसेच
  $\Sat{M}{\Gamma_2}$ आणि
  $\Gamma_2 \Entails \psi$ असल्यामुळे
  $\Sat{M}{\psi}$. या दोन्हींवरून
  $\Sat{M}{\varphi \land \psi}$.
\item शेवटचे अनुमान \Elim{\lor} आहे: सराव.
\item शेवटचे अनुमान \Elim{\lif} आहे. $\varphi \lif \psi$ आणि $\varphi$ या
  आधारविधानांपासून $\psi$ निष्पन्न केले जाते. निष्पत्ती~$\delta$
  पुढीलप्रमाणे दिसते:
  \begin{prooftree}
    \AxiomC{$\Gamma_1$}
    \RightLabel{$\delta_1$}
    \DeduceC{$\varphi \lif \psi$}
    \AxiomC{$\Gamma_2$}
    \RightLabel{$\delta_2$}
    \DeduceC{$\varphi$}
    \RightLabel{\Elim{\lif}}
    \BinaryInfC{$\psi$}
  \end{prooftree}
  विगमन गृहीतकानुसार $\varphi \lif \psi$ हे $\delta_1$ च्या
  मुक्त न केलेली गृहीतकांपासून~$\Gamma_1$ निष्पन्न होते आणि $\varphi$ हे
  $\delta_2$ च्या मुक्त न केलेली गृहीतकांपासून~$\Gamma_2$ निष्पन्न होते.
  एखादी
  संरचना~$\Struct{M}$
  विचारात घ्या. जर $\Sat{M}{\Gamma_1 \cup
    \Gamma_2}$, तर $\Sat{M}{\psi}$, हे दाखवायचे आहे.
  $\Sat{M}{\Gamma_1 \cup \Gamma_2}$ असे समजा.
  $\Gamma_1 \Entails \varphi \lif \psi$ असल्यामुळे
  $\Sat{M}{\varphi \lif \psi}$. तसेच
  $\Gamma_2 \Entails \varphi$ असल्यामुळे
  $\Sat{M}{\varphi}$. याचा अर्थ
  $\Sat{M}{\psi}$ असा होतो (कारण
  $\Sat/{M}{\psi}$ असेल, तर
  $\Sat{M}{\varphi}$ असल्यामुळे
  $\Sat/{M}{\varphi \lif \psi}$ असे होईल, जे
  $\Sat{M}{\varphi \lif \psi}$ च्या विरुद्ध आहे).
\item शेवटचे अनुमान \Elim{\lnot} आहे: सराव.
\item शेवटचे अनुमान \Elim{\lexists} आहे: सराव.
\end{enumerate}
\end{proof}
\begin{prob}
\ref{fol:ntd:sou:thm:soundness} ची सिद्धता पूर्ण करा.
\end{prob}
\begin{cor}
\label{fol:ntd:sou:cor:weak-soundness}
जर $\Proves \varphi$, तर $\varphi$ हे वैध आहे.
\end{cor}
\begin{cor}
\label{fol:ntd:sou:cor:consistency-soundness}
जर $\Gamma$ पूर्ततायोग्य असेल, तर तो सुसंगत आहे.
\end{cor}
\begin{proof}
प्रतिलोम विधान सिद्ध करू. $\Gamma$ सुसंगत नाही असे समजा. मग
$\Gamma \Proves \lfalse$, म्हणजे $\Gamma$ मधील मुक्त न केलेली
गृहीतकांपासून $\lfalse$ ची निष्पत्ती आहे. \ref{fol:ntd:sou:thm:soundness}
नुसार $\Gamma$ पूर्ण करणारी कोणतीही
संरचना~$\Struct{M}$
$\lfalse$ देखील पूर्ण करते. प्रत्येक
संरचना~$\Struct{M}$
साठी $\Sat/{M}{\lfalse}$ असल्यामुळे कोणतीही
$\Struct{M}$ $\Gamma$ पूर्ण करू शकत नाही,
म्हणजे $\Gamma$ पूर्ततायोग्य नाही.
\end{proof}
\section{निष्पत्ती सह एकरूपता}\label{fol:ntd:ide:sec}
एकरूपता असलेल्या निष्पत्त्या साठी अतिरिक्त अनुमाननियम आवश्यक असतात.
\par\medskip\noindent
\begin{defish}
\AxiomC{}
\RightLabel{\Intro{\eq}}
\UnaryInfC{$\eq[t][t]$}
\DisplayProof
\par\medskip\noindent
\begin{tabular}{r}
\AxiomC{$\eq[t_1][t_2]$}
\AxiomC{$\varphi(t_1)$}
\RightLabel{\Elim{\eq}}
\BinaryInfC{$\varphi(t_2)$}
\DisplayProof
\\[3ex]
\AxiomC{$\eq[t_1][t_2]$}
\AxiomC{$\varphi(t_2)$}
\RightLabel{\Elim{\eq}}
\BinaryInfC{$\varphi(t_1)$}
\DisplayProof
\end{tabular}
\end{defish}
वरील नियमांमध्ये $t$, $t_1$ आणि $t_2$ ही बंद पदे आहेत. \Intro{\eq} नियम
आपल्याला कोणत्याही गृहीतकांशिवाय, थेट $\eq[t][t]$ या रूपाचे कोणतेही
एकरूपता-विधान निष्पन्न करू देतो.
\begin{ex}
$s$ आणि $t$ ही बंद पदे असतील, तर $\varphi(s), \eq[s][t] \Proves \varphi(t)$:
\begin{prooftree}
\AxiomC{$\eq[s][t]$}
\AxiomC{$\varphi(s)$}
\RightLabel{$\Elim{\eq}$}
\BinaryInfC{$\varphi(t)$}
\end{prooftree}
याला “एकरूप वस्तूंच्या प्रतिस्थापनाचे तत्त्व” किंवा लाइब्निझचा नियम
म्हणून ओळखले जाऊ शकते.
\end{ex}
\begin{prob}
$=$ हे सममित आणि संक्रमक दोन्ही आहे हे सिद्ध करा; म्हणजे
$\lforall[x][\lforall[y][(\eq[x][y] \lif
    \eq[y][x])]]$ आणि $\lforall[x][\lforall[y][\lforall[z]((\eq[x][y]
    \land \eq[y][z]) \lif \eq[x][z])]]$ यांच्या
निष्पत्त्या द्या.
\end{prob}
\begin{ex}
आपण पुढील वाक्य निष्पन्न करतो:
\begin{align*}
& \lforall[x][\lforall[y][((\varphi(x) \land \varphi(y)) \lif \eq[x][y])]]
\intertext{या वाक्य पासून:}
& \lexists[x][\lforall[y][(\varphi(y) \lif \eq[y][x])]]
\end{align*}
निष्पत्ती उलट्या दिशेने विकसित करू:
\begin{prooftree}
\AxiomC{$\lexists[x][\lforall[y][(\varphi(y) \lif \eq[y][x])]]
\quad \Discharge{\varphi(a) \land \varphi(b)}{1}$}
\DeduceC{$\eq[a][b]$}
\DischargeRule{\Intro{\lif}}{1}
\UnaryInfC{$((\varphi(a) \land \varphi(b)) \lif \eq[a][b])$}
\RightLabel{\Intro{\lforall}}
\UnaryInfC{$\lforall[y][((\varphi(a) \land \varphi(y)) \lif \eq[a][y])]$}
\RightLabel{\Intro{\lforall}}
\UnaryInfC{$\lforall[x][\lforall[y][((\varphi(x) \land \varphi(y)) \lif \eq[x][y])]]$}
\end{prooftree}
आता मुख्य गृहीतक वापरावे लागेल: ते अस्तित्ववाची सूत्र असल्यामुळे
\Elim{\lexists} वापरून मध्यवर्ती निष्कर्ष~$\eq[a][b]$ निष्पन्न करू.
\begin{prooftree}
\AxiomC{$\lexists[x][\lforall[y][(\varphi(y) \lif \eq[y][x])]]$}
\AxiomC{$\Discharge{\lforall[y][(\varphi(y) \lif \eq[y][c]])}{2}$}
\noLine
\UnaryInfC{$\Discharge{\varphi(a) \land \varphi(b)}{1}$}
\DeduceC{$\eq[a][b]$}
\DischargeRule{\Elim{\lexists}}{2}
\BinaryInfC{$\eq[a][b]$}
\DischargeRule{\Intro{\lif}}{1}
\UnaryInfC{$((\varphi(a) \land \varphi(b)) \lif \eq[a][b])$}
\RightLabel{\Intro{\lforall}}
\UnaryInfC{$\lforall[y][((\varphi(a) \land \varphi(y)) \lif \eq[a][y])]$}
\RightLabel{\Intro{\lforall}}
\UnaryInfC{$\lforall[x][\lforall[y][((\varphi(x) \land \varphi(y)) \lif \eq[x][y])]]$}
\end{prooftree}
वरच्या उजव्या बाजूची उप-निष्पत्ती तिच्या गृहीतकांचा वापर करून
$\eq[a][c]$ आणि $\eq[b][c]$ दाखवल्यावर पूर्ण होते. यासाठी दोन स्वतंत्र
निष्पत्त्या आवश्यक आहेत. $\eq[a][c]$ साठीची निष्पत्ती
पुढीलप्रमाणे आहे:
\begin{prooftree}
\AxiomC{$\Discharge{\lforall[y][(\varphi(y) \lif \eq[y][c]])}{2}$}
\RightLabel{\Elim{\lforall}}
\UnaryInfC{$\varphi(a) \lif \eq[a][c]$}
\AxiomC{$\Discharge{\varphi(a) \land \varphi(b)}{1}$}
\RightLabel{\Elim{\land}}
\UnaryInfC{$\varphi(a)$}
\RightLabel{\Elim{\lif}}
\BinaryInfC{$\eq[a][c]$}
\end{prooftree}
$\eq[a][c]$ आणि $\eq[b][c]$ यांच्यापासून \Elim{\eq} वापरून
$\eq[a][b]$ निष्पन्न करतो.
\end{ex}
\begin{prob}
पुढील सूत्रांच्या निष्पत्त्या द्या:
\begin{enumerate}
\item $\lforall[x][\lforall[y][((\eq[x][y] \land \varphi(x)) \lif \varphi(y))]]$
\item $\lexists[x][\varphi(x)] \land \lforall[y][\lforall[z][((\varphi(y) \land
    \varphi(z)) \lif \eq[y][z])]] \lif \lexists[x][(\varphi(x) \land
  \lforall[y][(\varphi(y) \lif \eq[y][x])])]$
\end{enumerate}
\end{prob}
\section{एकरूपता सह निर्दोषता}\label{fol:ntd:sid:sec}
\begin{prop}
$\eq$ साठीच्या नियमांसह नैसर्गिक निगमन निर्दोष आहे.
\end{prop}
\begin{proof}
$\eq[t][t]$ या रूपातील कोणतेही सूत्र वैध आहे, कारण प्रत्येक
संरचना~$\Struct M$ साठी $\Sat{M}{\eq[t][t]}$. (लक्षात घ्या की
$t$ हे बंद पद आहे असे आपण गृहीत धरतो, म्हणजे त्यात कोणतेही चल नाही;
म्हणून चल-विन्यास अप्रासंगिक असतात.)
निष्पत्तीमधील शेवटचे अनुमान \Elim{\eq} आहे असे समजा, म्हणजे
निष्पत्तीचे रूप पुढीलप्रमाणे आहे:
\begin{prooftree}
  \AxiomC{$\Gamma_1$}
  \RightLabel{$\delta_1$}
  \DeduceC{$\eq[t_1][t_2]$}
  \AxiomC{$\Gamma_2$}
  \RightLabel{$\delta_2$}
  \DeduceC{$\varphi(t_1)$}
  \RightLabel{\Elim{\eq}}
  \BinaryInfC{$\varphi(t_2)$}
\end{prooftree}
आधारविधाने~$\eq[t_1][t_2]$ आणि $\varphi(t_1)$ ही अनुक्रमे
मुक्त न केलेली गृहीतके~$\Gamma_1$ आणि $\Gamma_2$ यांच्यापासून
निष्पन्न केली आहेत. $\varphi(t_2)$ हे $\Gamma_1 \cup \Gamma_2$ पासून
निष्पन्न होते, हे दाखवायचे आहे. $\Sat{M}{\Gamma_1 \cup
  \Gamma_2}$ असलेली संरचना~$\Struct{M}$ विचारात घ्या. विगमन
गृहीतकानुसार $\Sat{M}{\varphi(t_1)}$ आणि $\Sat{M}{\eq[t_1][t_2]}$.
त्यामुळे $\Value{t_1}{M} = \Value{t_2}{M}$. $s$ हा कोणताही चल-विन्यास
आणि $m = \Value{t_1}{M} = \Value{t_2}{M}$ असू द्या.
\ref{fol:syn:ext:prop:ext-formulas} नुसार
$\Sat{M}{\varphi(t_1)}[s]$ तेव्हा आणि केवळ तेव्हाच, जेव्हा
$\Sat{M}{\varphi(x)}[\Subst{s}{m}{x}]$, आणि हे तेव्हा आणि केवळ तेव्हाच, जेव्हा
$\Sat{M}{\varphi(t_2)}[s]$. $\Sat{M}{\varphi(t_1)}$ असल्यामुळे
$\Sat{M}{\varphi(t_2)}$.
\end{proof}
